\documentclass[aps,pra,twocolumn,superscriptaddress,longbibliography]{revtex4-2}

\usepackage{amsmath,amssymb,amsthm,graphicx,xcolor,bm,braket,dcolumn,booktabs}
\usepackage{hyperref}

\newtheorem{theorem}{Theorem}
\newtheorem{lemma}{Lemma}
\newtheorem{proposition}{Proposition}
\newtheorem{conjecture}{Conjecture}
\newtheorem{observation}{Observation}

% Numeric-registry citation macro (CLAUDE.md SS15): renders the literal
% only; the key is a foreign key into debug/qa/numeric_registry.py.
\newcommand{\gvq}[2]{#2}

\begin{document}

\title{Linked Electron Densities Climb the Banana Ladder:\\
Period Classes of Molecular Integrals over Exponential-Type Orbitals}
\thanks{Paper 63 in the Geometric Vacuum series}

\author{J.~Loutey}
\affiliation{Independent Researcher, Kent, Washington}

\date{November 6, 2026}

%% ========================================================================
%% ABSTRACT
%% ========================================================================

\begin{abstract}
Paper~59 read the radial core of the three-center electron-repulsion
integral over $1s$ Slater orbitals as an elliptic Bessel moment:\ two
two-center densities coupled by a Coulomb kernel build a genus-one curve.
This paper asks two questions that result leaves open.  How does the class
climb when an integral links more densities?  And how far does the reading
reach past $1s$ orbitals?  \textbf{[SYMBOLIC PROOF + PRIOR ART]} The
second is answered first, by a dictionary.  Every exponential-type orbital
on a center is a polynomial differential operator in its exponent and its
center coordinate, with coefficients rational in the exponent, applied to
the $1s$ Yukawa function on that center;\ at the fibre, for $s$-type
densities, every density weight is a polynomial in the Euler operators
$c\,\partial_c$ applied to the $1s$ weight.  So every multi-density integral with rational kernels is that
operator applied to the $1s$ integral of the same density graph, its
$D$-module is a quotient of the $1s$ one, and its rung is at most the $1s$
rung of its graph.  The ladder below is therefore a property of the graph,
not of the orbitals on it.  \textbf{[SYMBOLIC + MEASURED]} On the graph the
rung is read off the holomorphic period.  Any tree of $n$ linked
two-center densities has the same period, a Bessel moment in $n$ scales;\
its one-scale family is annihilated by an operator of order $n$ that is the
Hadamard product of the Legendre operator with the operator of the equal-mass
banana integral with $n-1$ propagators:\ $n=2$ is Legendre's $K$, $n=3$
Clausen's elliptic square (a K3 surface), $n=4$ entry \#58 of the
Calabi--Yau tables, the first period not built from elliptic ones, and $n=5$
brings no new type.  The real three-dimensional star and path of four
densities reduce onto that operator with no K3 piece.  A fourth distinct
center adds no genus \textbf{[MEASURED]}.  \textbf{[SYMBOLIC + MEASURED]} A
Coulomb ring of three densities does not climb:\ its loop integrates out by
the three-dimensional uniqueness relation and it sits on the elliptic rung,
below the chain of three;\ the correlated triangle of explicitly correlated
methods, with exponential geminals on two edges, sits with the chain on the
K3 rung, its geminal acting as an imaginary external momentum, and its
operators follow from integration by parts on the kite and the tetrahedron.
A Gaussian geminal leaves the ladder for exponential periods.
\textbf{[OBSERVATION]} Read as a Feynman graph, every instance fits one
pattern:\ the class is set by the shape of the density graph and the masses
on its lines.  We state it as a bounded conjecture.  Trees sit at rung
$n$;\ the rung names an operator and not a finite set of numbers, since on
every slice tested the three-electron fibre is a K3 function integrated
once and at generic scales it has no relation to the K3 and elliptic periods
in any basis tried;\ rings are ordered by uniqueness, not by count;\ and the
integrated value sits at or below its fibre.  The caps are stated:\ the
three-dimensional reduction is shown through four densities, the only ring
tested is the triangle, and the generic fibre's content is evidence.
\textbf{[PRIOR ART]} The operators are known in the mathematics of Bessel
moments and banana integrals;\ the Averys' momentum-space treatment of these
same integrals re-expands every two-center density in a one-center basis and
reaches the genus-zero class only through an infinite, non-sparse series.
What is new is the route from molecular integrals to the ladder, the
dictionary that makes the ladder orbital-independent, and the prospective
record on which the rungs were called before they were computed.
\end{abstract}

\maketitle

%% ========================================================================
\section{Introduction}
\label{sec:intro}
%% ========================================================================

Two-center electron-repulsion integrals over exponential-type orbitals close
in elementary functions together with the exponential integral, the
logarithm and Euler's constant:\ a genus-zero class, weight one, settled by
Roothaan and Ruedenberg.  Paper~59~\cite{loutey_paper59} read the genuine
three-center integral $(XY|XZ)$ in momentum space, where each center is a
phase and the two two-center densities carry Fock dispersion roots
$\Delta_i=\sqrt{c_ik^2+1}$ with different scales, and found its radial core
on the elliptic curve
%
\begin{equation}
  y^{2} = (c_{1}k^{2}+1)(c_{2}k^{2}+1),
  \label{eq:curve}
\end{equation}
%
whose period at zero separation is a complete elliptic integral,
%
\begin{equation}
  \int_{0}^{\infty}\!\frac{dk}{\sqrt{(c_{1}k^{2}+1)(c_{2}k^{2}+1)}}
  = \frac{1}{\sqrt{c_{\max}}}\,K(m),\qquad m=1-\frac{c_{\min}}{c_{\max}},
  \label{eq:period}
\end{equation}
%
degenerating to the elementary $\pi/(2\sqrt c)$ only at equal scales.  That
paper's object was one integral.  This paper's object is the \emph{ladder}:\
what happens to the class as an integral links more electron densities, and
whether the reading depends on the $1s$ orbitals it was found on.

The answer to the second question is a dictionary
(Sec.~\ref{sec:dictionary}), and it is classical in its parts.  Every
exponential-type orbital is a derivative of the $1s$ Yukawa function in its
exponent and its center;\ so every integral over such orbitals is a
derivative of the $1s$ integral of the same \emph{density graph}, the graph
whose lines are the orbitals and the kernels and whose vertices are the
electrons and the nuclei.  Derivatives in the parameters of a period stay in
its $D$-module.  The rung is therefore a property of the graph.  The answer
to the first question (Sec.~\ref{sec:beyond}) was found on $1s$ orbitals and,
by the dictionary, holds for all of them:\ for a tree of $n$ linked densities
the holomorphic period depends only on $n$, and the operators of the one-scale
family for $n=2,\dots,5$ are, under one exact Hadamard product, the operators
of the equal-mass banana Feynman integrals, which carry Calabi--Yau
$(n-2)$-folds~\cite{bonisch2021}.  The Coulomb ring of three does not climb,
the correlated triangle does, and a Gaussian geminal leaves the ladder.

The reading that organizes these facts is a Feynman-graph one
(Sec.~\ref{sec:conjecture}):\ densities and kernels are lines, exponents and
geminal decays are masses, and the class is set by the shape of the graph and
the masses on its lines.  We state it as a bounded conjecture with its
falsifiers, and we record the one place its naive form fails:\ the rung names
an operator, not a finite basis of periods.  In the transcendence bookkeeping
of the Geometric Vacuum program (Paper~18~\cite{loutey_paper18}) this gives the
many-electron integrals a second coordinate beside the Tate weight of the
one-body and spectral projections:\ the Calabi--Yau dimension of the density
graph's banana.  The seam catalogue (Paper~57~\cite{loutey_paper57}) records
the rows;\ the number-theoretic structure of the first rung is the companion
Paper~61~\cite{loutey_paper61}.

Most of the mathematics below is prior art, and Sec.~\ref{sec:placement}
says whose.  The route from a molecular integral to it was not found in a
targeted search, and the rungs were called in committed records before they
were computed (Sec.~\ref{sec:beyond} names each record).  The caps are
stated in Sec.~\ref{sec:scope}.

%% ========================================================================
\section{The dictionary:\ every exponential-orbital integral is a derivative of a $1s$ Yukawa integral}
\label{sec:dictionary}
%% ========================================================================

\emph{Setting.}  An exponential-type orbital on a center $\bm A$ is a finite
sum of terms $|\bm r-\bm A|^{j}\,e^{-\zeta|\bm r-\bm A|}\,\mathcal{Y}_{\ell m}(\bm r-\bm A)$
with $j\ge0$ an integer and $\mathcal{Y}_{\ell m}$ a regular solid harmonic
(homogeneous of degree $\ell$).  Slater orbitals are the single terms
$j=n-1-\ell$;\ hydrogenic and Coulomb--Sturmian orbitals, a Laguerre
polynomial times an exponential, are such finite sums.  The $1s$ seed is the
Yukawa function $g_\zeta(s)=e^{-\zeta s}/s$, $s=|\bm r-\bm A|$, the
three-dimensional Green's function of $-\nabla^2+\zeta^2$, whose Fourier
transform is $4\pi/(p^2+\zeta^2)$.  A \emph{density graph} $\Gamma$ has a
vertex for every electron (integrated) and every nucleus (fixed), a line of
mass $\zeta$ from a nucleus to an electron for every orbital, and a line
between two electrons for every kernel:\ massless for the Coulomb kernel,
mass $\gamma$ for a Slater geminal $e^{-\gamma r_{12}}$, which is itself
$-\partial_\gamma$ of the Yukawa line.  The \emph{$1s$ integral of $\Gamma$},
$I^{(1s)}_\Gamma(\zeta_e;\bm X_v)$, is the integral with every orbital line
the Yukawa $g_{\zeta_e}$.

\begin{lemma}[orbital level; prior art]
\label{lem:orbital}
Every exponential-type orbital on $\bm A$ is
$\mathcal{D}(\partial_\zeta,\nabla_{\!\bm A})\,g_\zeta(|\bm r-\bm A|)$ for a
polynomial differential operator $\mathcal{D}$ whose coefficients are
rational in $\zeta$.
\end{lemma}

\textbf{[SYMBOLIC PROOF]} Two identities do it.  Powers are
$\zeta$-derivatives, $s^{j}e^{-\zeta s}=(-\partial_\zeta)^{j+1}g_\zeta(s)$.
Solid harmonics are center derivatives:\ by Hobson's theorem,
$\mathcal{Y}_{\ell m}(\nabla)F(s)=\mathcal{Y}_{\ell m}(\bm r-\bm A)\,(s^{-1}d/ds)^{\ell}F(s)$
for any radial $F$, and $\nabla=-\nabla_{\!\bm A}$ on a function of
$\bm r-\bm A$.  Applied to $F=s^{m-1}e^{-\zeta s}$ this gives
$\mathcal{Y}_{\ell m}(\bm r-\bm A)\,e^{-\zeta s}\sum_{i=0}^{\ell}a_i(m,\ell)\,\zeta^{\,i}\,s^{m-1-2\ell+i}$
with $a_\ell=(-1)^\ell$;\ the system in $m$ is triangular in the power of $s$
with leading coefficient $(-\zeta)^{\ell}\neq0$, so every
$\mathcal{Y}_{\ell m}(\bm r-\bm A)\,s^{j}e^{-\zeta s}$ is a combination of the
$\mathcal{Y}_{\ell m}(-\nabla_{\!\bm A})(-\partial_\zeta)^{m}g_\zeta$ with
coefficients in $\mathbb{Q}[\zeta,\zeta^{-1}]$.  (For $\ell=1$:\
$\mathcal{Y}_{1m}e^{-\zeta s}=-\zeta^{-1}[\,T_2+\zeta^{-1}T_1\,]$ with
$T_m=\mathcal{Y}_{1m}(-\nabla_{\!\bm A})(-\partial_\zeta)^{m}g_\zeta$.)
In momentum space the same statement reads:\ the transform of every
exponential-type orbital is $e^{-i\bm p\cdot\bm A}\,\mathcal{Y}_{\ell m}(i\bm p)\,R(p^2,\zeta)$
with $R$ rational, since $\mathcal{Y}_{\ell m}(i\nabla_{\!\bm p})$ acting on a
rational function of $p^2$ gives $\mathcal{Y}_{\ell m}(i\bm p)$ times a
rational function.  \textbf{[PRIOR ART]} Both halves are classical:\ the
rational momentum-space form of exponential-type orbitals is Steinborn and
Filter's~\cite{weniger2005}, generating angular momentum by differentiating
in the nuclear coordinates is standard~\cite{weniger2005}, and the
$\zeta$-derivative is the device Paper~59 already uses to restore a Slater
orbital from a Yukawa.  Guard:\ \texttt{tests/test\_paper63\_dictionary.py}
(the $\ell=1$ inversion and the $\ell=2$ triangular system in exact
arithmetic;\ the rational momentum-space form for $(n,\ell)$ up to $(4,3)$).

\begin{lemma}[integral level]
\label{lem:integral}
Let $I_\Gamma$ be a multi-density integral over exponential-type orbitals
with Coulomb or Slater-geminal kernels.  Then
$I_\Gamma=\prod_e\mathcal{D}_e\;I^{(1s)}_\Gamma$, the operators acting on the
exponents and centers of the $1s$ integral of the same graph.
\end{lemma}

\textbf{[SYMBOLIC PROOF]} Insert Lemma~\ref{lem:orbital} for each orbital.
The $1s$ integrand is bounded by a product of exponentials in the
$|\bm r_i-\bm X_v|$ times integrable Coulomb singularities, uniformly for
exponents in a compact subset of $(0,\infty)$ and centers in a compact set, as
are its derivatives of every order in those parameters;\ dominated convergence
moves the $\mathcal{D}_e$ outside the integral.  The same holds for the
momentum form, where the orbital lines are rational and the exponential decay
of position space becomes the polynomial decay of the integrand in momentum.

\begin{proposition}[the rung is a property of the graph]
\label{prop:rung}
Let $\mathcal{M}$ be the holonomic $D$-module generated by $I^{(1s)}_\Gamma$
over the exponents and centers.  Every $I_\Gamma$ over exponential-type
orbitals lies in $\mathcal{M}$;\ the composition factors of its minimal
annihilator in any one parameter are among those of $I^{(1s)}_\Gamma$'s.  At
the fibre, where the centers coincide, a density built from $s$-type orbitals
enters through its dispersion root $\Delta_v=\sqrt{c_vL_v^2+1}$ alone, and
every such weight $P(L_v^2)\,\Delta_v^{-(2m+1)}$ is a polynomial in $c_v$,
$c_v^{-1}$ and the Euler operator $c_v\partial_{c_v}$ applied to
$\Delta_v^{-1}$ (for $\ell>0$ the center derivatives of
Lemma~\ref{lem:orbital} leave polynomial momentum factors at the fibre, which
the first statement covers and these identities do not):
%
\begin{equation}
  \Delta^{-(2m+1)}=\Bigl(1+\tfrac{2}{2m-1}\,c\,\partial_c\Bigr)\Delta^{-(2m-1)},
  \qquad
  L^{2}\,\Delta^{-j}=c^{-1}\bigl(\Delta^{-(j-2)}-\Delta^{-j}\bigr).
  \label{eq:euler_weights}
\end{equation}
\end{proposition}

\textbf{[SYMBOLIC PROOF]} The first statement is
Lemma~\ref{lem:integral} read in $\mathcal{M}$:\ a section of a holonomic
module and its derivatives generate a submodule, so the annihilator of
$\mathcal{D}\,I^{(1s)}$ contains a left multiple of that of $I^{(1s)}$ and
its composition factors are a subset.  For the fibre identities,
$\partial_c\Delta^{-1}=-\tfrac12L^2\Delta^{-3}$ and $L^2=(\Delta^2-1)/c$ give
$c\,\partial_c\Delta^{-1}=-\tfrac12(\Delta^{-1}-\Delta^{-3})$, which is the
case $m=1$;\ induction on $m$ gives the rest, and the second identity is
$L^2=(\Delta^2-1)/c$ again.  \textbf{[MEASURED]} On the two-density rung the
identities are checked against the period~\eqref{eq:period}:\
$\int_0^\infty\Delta_1^{-3}\Delta_2^{-3}dk$ and
$\int_0^\infty\Delta_1^{-5}\Delta_2^{-5}dk$ equal the operators of
Eq.~\eqref{eq:euler_weights} applied to $K(m)/\sqrt{c_{\max}}$, to $10^{-25}$
at three scale pairs, and the result is again a combination of $K$ and $E$:\
the Legendre rung, as the proposition requires.  On the three-density rung
the Euler operators act under the tree identity (Eq.~\eqref{eq:tree} below):\
the one-dimensional transform of the $1s$ weight is $2c^{-1/2}K_0(x/\sqrt c)$,
and Eq.~\eqref{eq:euler_weights} turns it into $2xc^{-1}K_1(x/\sqrt c)$ for
$\Delta^{-3}$ and $\tfrac43xc^{-1}K_1+\tfrac23x^2c^{-3/2}K_0$ for
$\Delta^{-5}$, the known transforms of those powers.  The two-dimensional
fingerprint integral of the chain of three, with $\Delta^{-3}$ or
$\Delta^{-5}$ on every density, equals $\pi^{-1}\!\int_0^\infty\prod_v w_v(x)\,dx$
with those weights to $2\times10^{-9}$ at three scale triples, the precision of
the two-dimensional quadrature;\ the $1s$ weight in place of $w_v$ misses by
more than a percent.  Guard:\ \texttt{tests/test\_paper63\_dictionary.py}
(29 checks, fire-tested).

\emph{What the dictionary does and does not say.}  It lifts every statement
below from $1s$ orbitals to all exponential-type orbitals \emph{as an upper
bound}:\ a general-orbital integral cannot climb above the $1s$ integral of
its graph.  It does not say the rung is attained;\ a derivative can land in a
submodule.  In every case computed it is attained---Paper~59's chain at the
Slater weights has the same composition factors as the chain at the
fingerprint weights (Sec.~\ref{sec:density_ring})---and the proposition is
why that had to be so or lower, not a proof that it is equal.  Two caps it
inherits unchanged from Paper~59:\ the rung is read on the fibre and its
holomorphic period, and the integrated value, over the Feynman parameters
and separations, may sit lower;\ and the reduction of the real
three-dimensional integral to the one-dimensional Bessel-moment family, which
the dictionary does not supply, is shown below for the chain of three and the
two trees of four only.  Gaussian orbitals and Gaussian geminals are outside
the lemma:\ $e^{-\alpha r^2}$ is no derivative of a Yukawa, its transform is
not rational, and Sec.~\ref{sec:density_ring} finds its integrals off the
ladder.

%% ========================================================================
\section{The rungs:\ a fourth center, a three-electron chain, four linked densities, and a ring}
\label{sec:beyond}
%% ========================================================================

Two results pre-registered on 2026-10-13 (records
\texttt{docs/predictions/2026-10-13\_four\_centre\_genus.md} and
\texttt{docs/predictions/2026-10-13\_three\_electron\_k3.md}, each committed
before its computation) extend the genus reading of
Paper~59~\cite{loutey_paper59} in two directions:\ more centers, and more Coulomb
kernels.  A third, registered on 2026-10-16
(\texttt{docs/predictions/2026-10-16\_linking\_ceiling.md}, written before
its computation but not committed beforehand), asks how the rung grows with
the number of linked densities.  A fourth, registered and committed on
2026-10-17 before its computation
(\texttt{docs/predictions/2026-10-17\_density\_loop.md}), asks whether a
ring of linked densities climbs faster than a tree.  A fifth, committed the
same day before its computation
(\texttt{docs/predictions/2026-10-17\_ring\_extensions.md}), asks whether the
correlated triangle, with geminals on two edges, behaves like the Coulomb
one, and how a ring of four behaves;\ the second test stopped for feasibility.
A sixth, committed before its computation
(\texttt{docs/predictions/2026-10-17\_ising\_class.md}), asks where the
Ising-class moments sit on the ladder.

\subsection{A fourth distinct center adds no genus}
\label{sec:four_centre}

\textbf{[MEASURED]} The general four-center integral $(XY|ZW)$, four
distinct centers and $1s$ Slater orbitals, reduces by the same Feynman
parametrization to one radial $k$-integral over two parameters $(s,t)$
whose integrand carries exactly the two dispersion roots
$\Delta_1(s)$, $\Delta_2(t)$ of the curve~\eqref{eq:curve};\ the centers enter
only through a single phase vector $\bm W(s,t)$ and the two pair
distances $D_1$, $D_2$.  The reduced formula gives $0.218814149639$
against $0.218814149425$ from McMurchie--Davidson on sixteen-Gaussian
fits (the residual is the fit gap), and reproduces the three-center
value $0.204941652531$ of Paper~59 as a special case.
There is no genus two at four centers, consistent with the class being
set by the number of distinct two-center densities the kernel couples,
not by the number of nuclei;\ this is that rule's first instance beyond
the one this paper was written on.  Its honest weight is that of a
corollary of Paper~59's reduction, a calibration rather than a hard
test.  The momentum form itself is the Averys'
(Ref.~\cite{avery_avery_4center}, App.~B);\ the genus reading is not in
it.  Guards:\ \texttt{tests/test\_paper63\_four\_centre\_genus.py} (both
values recomputed against McMurchie--Davidson;\ the kernel derived
symbolically as the second $z$-derivative of a Yukawa product;\ the
integrand's square roots and the dependence on the pair placement
checked).

\subsection{A three-electron chain lives on a K3 surface}
\label{sec:three_electron}

\textbf{[SYMBOLIC + MEASURED]} The chain integral
$\langle \rho_a(1)\rho_b(2)\rho_c(3)/(r_{12}r_{13})\rangle$, three
two-center $1s$ Slater transition densities (dispersion roots
$\Delta_b = \sqrt{c_b p^2 + 1}$, $\Delta_c = \sqrt{c_c q^2 + 1}$,
$\Delta_a = \sqrt{c_a |p+q|^2 + 1}$, scales $a = c_a^{-1/2}$, $b = c_b^{-1/2}$,
$c = c_c^{-1/2}$) joined by two Coulomb kernels at electron~1, has two momentum integrations and three
dispersion roots.  The angular integral is elementary (no logarithm);\
the reduced two-dimensional fibre matches the direct three-dimensional
one ($10^{-15}$ to $4\times10^{-14}$ in the record, $10^{-9}$ in the
reduced-precision guard);\ the one-center limit is rational
($23/54$, $133/144$, $289/225$ at the three registered exponent triples);\
the branch locus in $\mathbb{P}^2$ is six lines with no four concurrent,
so the double cover is a K3 surface, and fixed-$p$ slices are
non-isotrivial genus-one curves.  \textbf{[MEASURED / classical]} The
holomorphic period reduces by one Fourier transform to the three-scale
Bessel moment $\bar V_3 = \int_0^\infty K_0(at)K_0(bt)K_0(ct)\,dt$ (the
notation of Ref.~\cite{bbbg2008}), whose closed form for $c \ge b \ge a > 0$
is $(2c/\pi)\,\bar V_3 = K(k_-)K(k_+') + K(k_+)K(k_-')$, with the moduli
$k_\pm$ of Ref.~\cite{bbbg2008} (Eq.~33, obtained there from an identity of
Bailey~\cite{bailey1936}), checked at four scale triples to $10^{-34}$ in
the record and at three in the guard to $10^{-17}$.
The K3's period is therefore ``elliptic squared'':\ it opens no new
transcendence class, being a sum of products of periods of the genus-one class of Paper~59.
The reduction from a Slater chain to $\bar V_3$ is this paper's;\ the closed
form is that of Ref.~\cite{bbbg2008}.  In the language of Paper~59's literature placement, the
repulsion integral sits on the rung of the two-loop sunrise and the chain
on the rung of the three-loop banana;\ the families differ, and
Sec.~\ref{sec:four_density} gives the exact relation.  \textbf{[OBSERVATION, evidence]} At equal scales the
true fibre at zero separation collapses to two K3 periods (Bessel moments
of the record) and $\pi^3$.  At four generic scale triples, $(1;2,4)$,
$(1;3,7)$, $(3;1,7)$ and $(3;4,6)$ with the central density's scale first,
integer-relation searches at 340 and 400 digits find no relation in five
bases built from the K3 periods and their scale derivatives, products of
elliptic periods at the Bailey moduli, residue periods and their
products, dilogarithms and trilogarithms of scale ratios, the vacuum
sunrise and its mass derivatives, logarithms times periods, and
weight-three Tate numbers.  Every basis recovers a planted relation, and
in every basis the fibre with $\int_0^x K_0$ dropped from each leaf
reduces onto the K3 periods;\ three of the twenty searches are void, their
precision cap below the planted relation's height.  Moving the central
density within the scale set $\{1,3,7\}$ gives no relation either, and the
fibre does not collapse at the two threshold triples $(2;1,1)$ and
$(3;1,2)$, where the central scale is the sum of the others.  That is
evidence, at a sample of triples and with guessed bases, that the generic
fibre carries content beyond the K3 and elliptic periods, as the banana
does;\ on a one-parameter slice it is a K3 function integrated once
(Sec.~\ref{sec:density_ring}), a class no finite basis of periods is
expected to capture.  \textbf{[SYMBOLIC + MEASURED]} A second slice, with
scales $(t;1,2)$, passes through two of these triples, $(1;2,4)$ at $t=1/2$
and the threshold triple $(3;1,2)$ at $t=3$, where the fibre values
agree with the slice to 20 digits.  Its Mellin operator has order 17 and
divides exactly over the factor $\theta+2$ and two first-order pieces,
one for each leaf, which enter together as their least common left
multiple.  Above them sits a piece of order four, equivalent by an exact
map with poles only at $t=0,\pm1,\pm3$ to the operator of the three-scale
moment $\int_0^\infty K_0(tx)K_0(x)K_0(2x)\,dx$.  The composite of order
seven annihilates the fibre to $5\times10^{-33}$ of its largest term at
32 digits, the fibre needs the order-four piece, and a search modulo a
prime finds no rational map splitting it off within the bound set by the
local exponents, while the same call finds a planted split.  So on this
slice too the fibre is an elementary function plus a K3 function
integrated once, which accounts for the null at those two triples
(record \texttt{docs/predictions/2026-10-21\_generic\_slice.md}).
Longer trees are taken up in Sec.~\ref{sec:four_density}.
Guards:\ \texttt{tests/test\_paper63\_three\_electron\_k3.py}
(the fibre agreement, the elementary angle integral, the rational
one-center limits, the closed form against quadrature, the K3 period
against $\bar V_3$, the six-line branch locus, the non-isotrivial slice);\
\texttt{tests/test\_paper63\_generic\_slice.py} (the second slice);\
the equal-scale collapse and the generic null are not guarded, and the
records carry them (\texttt{docs/predictions/2026-10-13\_three\_electron\_k3.md},
\texttt{docs/predictions/2026-10-19\_p3e2\_widened.md}).

\subsection{Four linked densities:\ the first rung beyond elliptic}
\label{sec:four_density}

\textbf{[SYMBOLIC + MEASURED]} The holomorphic period of a tree of linked
densities depends only on how many densities the kernels link, not on the
pattern of the links;\ a Coulomb ring of three behaves differently
(Sec.~\ref{sec:density_ring}).  Take $n$ two-center $1s$ densities joined into a tree by $n-1$
rational kernels, and write the dispersion root of density $v$ at the
fingerprint fibre as $\Delta_v=\sqrt{c_vL_v^2+1}$, with $L_v$ the signed
sum of the momenta on the kernels at $v$.  Fourier-transforming each
$1/\Delta_v$ to $K_0(|y|/\sqrt{c_v})/(\pi\sqrt{c_v})$ turns every momentum
integration into a delta function that identifies the two ends of its
kernel, and a tree is connected, so
\begin{equation}
  \int_{\mathbb{R}^{n-1}} d^{n-1}k\,\prod_{v=1}^{n}\frac{1}{\Delta_v(L_v(k))}
  = \frac{2^{n}}{\pi\sqrt{c_1\cdots c_n}}
    \int_0^\infty \prod_{v=1}^{n} K_0(a_v x)\,dx,
  \qquad a_v = c_v^{-1/2},
  \label{eq:tree}
\end{equation}
for a star and a chain alike.  At $n=3$ this is the $\bar V_3$ reduction
of Sec.~\ref{sec:three_electron};\ at $n=2$ it is the elliptic period of
Eq.~\eqref{eq:period}.  A three-dimensional evaluation of four-density
stars and chains, with the faster-decaying roots $\Delta_v^{-3}$ to which
the same argument applies, matches the Bessel-moment side to between
$4\times10^{-9}$ and $2\times10^{-7}$, the grid's own convergence.  That
check tests the identity's bookkeeping;\ the real three-dimensional
integral reduces to moments of the same family with other weights in $x$,
as shown below for both trees of four densities, the star and the path.

\textbf{[SYMBOLIC]} To read the rung, vary one scale:\
$F_n(t)=\int_0^\infty K_0(x)^{n-1}K_0(tx)\,dx$.  The Mellin transform of
$K_0^{\,n-1}$ obeys the recurrence of the $(n-1)$-th symmetric power of the
Bessel operator, and that of $K_0(tx)$ is $2^{s-2}\Gamma(s/2)^2$, so $F_n$
satisfies an explicit linear equation in $t$ of order $n$;\
\textbf{[MEASURED]} each one is checked against direct quadrature of $F_n$
to $10^{-28}$.  For $n=2$ it is the Legendre equation of $K$.  For $n=3$
it is the equation of ${}_3F_2(\tfrac12,\tfrac12,\tfrac12;1,1;t^2/4)$,
Clausen's square of an elliptic period, as Bailey's formula requires.
For $n=4$, in the variable $x=1/(4t^2)$, it is
\begin{equation}
  \theta^4 - x(2\theta+1)^2(10\theta^2+10\theta+3)
  + 9x^2(2\theta+1)^2(2\theta+3)^2,
  \qquad \theta = x\,\frac{d}{dx},
  \label{eq:four_density_op}
\end{equation}
which has maximally unipotent monodromy at $x=0$ and conifold points at
$x=1/36$ and $x=1/4$, with local exponents $(0,1,1,2)$.  Its holomorphic
solution is $\sum_m B_m (x/4)^m$ with integer
$B_m=\binom{2m}{m}^2\sum_k\binom{m}{k}^2\binom{2k}{k}$, the term-by-term
(Hadamard) product of $\binom{2m}{m}^2$ with the Ap\'ery-like sequence of
the equal-mass sunrise, and its mirror map has integer coefficients.

\textbf{[SYMBOLIC]} Whether a family is built from elliptic periods shows
in the quadratic relations among the solutions of its equation:\ for
monodromy of these types, an elliptic family has none, the square of
one elliptic family or the tensor product of two has one, and the cube of
one has three.  Counted
exactly on Taylor bases at two rational points, the numbers are
$0,1,0,1$ for $n=2,3,4,5$.  For $n=4$ the operator is irreducible, with
no bound on degrees:\ a right factor of order $k$ would carry, at each
singular point, $k$ of the operator's local exponents, and Fuchs' relation
then fails for every $k<4$, however many apparent singularities the factor
adds.  The same count shows the operators for $n=2,3,5$ irreducible.  Its
Yukawa coupling,
$1+16q+1152q^2+99520q^3+\cdots$, is not constant, as it would be for a
symmetric cube.  An irreducible fourth-order local system with no invariant
quadric that is not a symmetric cube has differential Galois group $\mathrm{Sp}_4$ or $\mathrm{SL}_4$,
and neither is a quotient of a product of copies of $\mathrm{SL}_2$.  So
$F_4$ is not a polynomial in elliptic periods with algebraic moduli, and
neither is the generic four-scale moment, which restricts to it.  Two
linked densities give an elliptic period, three give a K3 whose period is
elliptic squared, and four give the first period of a new kind.  At $n=5$
the single relation marks an orthogonal structure of rank five, not the
fourth power of an elliptic period, which would carry six.  That is the
structure of the exterior square of a fourth-order Calabi--Yau operator,
and the tables of Ref.~\cite{aesz2005} print the operator of $F_5$
verbatim, attached to entry \#39, a different operator from the \#58 of
$n=4$.  So five linked densities give no new type of period:\ the type is
that of four, in a different family.

\textbf{[OBSERVATION]} Every electron-repulsion integral couples exactly
two densities, so the integrals of a molecular Hamiltonian stay on the
elliptic rung however large the molecule.  The class climbs only in
correlated integrals that link more densities, such as the four-electron
bridge integrals of explicitly correlated methods taken over two-center
densities.  Two caps apply.  These are statements about the fibre and
its holomorphic period, so the integrated value may sit lower, as the
two-center exchange class does (Paper~59, geminal section).  And the reduction of
the real three-dimensional integral into the family of
Eq.~\eqref{eq:tree} is shown only for the chain of three and the two trees
of four below;\ for five or more densities it is not shown.  \textbf{[SYMBOLIC +
MEASURED]} At the fingerprint fibre on the slice $(t;1,1,1)$ the star is
$(16/\pi^4)\,t^3\int_0^\infty\chi(x)^3K_0(tx)\,dx/x$, with
$\chi=\int_0^x sK_1(s)\,ds$ (Sec.~\ref{sec:density_ring}), checked against
a momentum-space evaluation to $2\times10^{-21}$.  Its Mellin operator, of
order 19, divides exactly with a piece of order four equivalent to the
operator~\eqref{eq:four_density_op} by an exact rational map.  An
order-eight right factor annihilates the star to $3\times10^{-45}$ of its
largest term at 45 digits;\ its composition factors are that piece, the
Legendre operator of $F_2$ and two of order one, with no K3 piece, and the
star needs the top piece.  So the star sits on the rung of four linked
densities and skips the K3 rung;\ minimality is bounded by a search modulo
a prime, not proved (record
\texttt{docs/predictions/2026-10-21\_four\_density\_star.md}).  Guards:\
\texttt{tests/test\_paper63\_four\_density\_rung.py},
\texttt{tests/test\_paper63\_four\_density\_star.py}.
\textbf{[SYMBOLIC + MEASURED]} The path, four densities in a chain,
reduces the same way.  On the slice $(1,t,1,1)$, with $t$ on an inner
density, it is
$(16/\pi^4)\,t^3\int_0^\infty\!\int_0^\infty\min(x,y)\,K_0(tx)\chi(x)\,K_0(y)\chi(y)\,dx\,dy$,
checked against a momentum-space evaluation to 20 digits.  Its holonomic
relations give an operator of order 74 with an exact right factor of
order 12, which annihilates the path to below $10^{-28}$ of its largest
term at 32 digits and $4\times10^{-56}$ at 60.  That factor splits exactly
into a piece equivalent to the operator~\eqref{eq:four_density_op}, three
pieces equivalent to the Legendre operator of $F_2$, and two of order
one, each equivalence by an exact rational map;\ there is no K3 piece,
and the path needs the top piece.  So both trees of four densities sit on
the same rung and skip the K3 rung;\ for the path, minimality is bounded
by a search modulo a prime up to $t$-degree 260 (record
\texttt{docs/predictions/2026-10-22\_four\_density\_path.md}).  Guard:\
\texttt{tests/test\_paper63\_four\_density\_path.py}.

\textbf{[PRIOR ART]} Most of this structure is known in the mathematics of
Bessel moments;\ what is new here is the route into it from a molecular
integral.  The operator~\eqref{eq:four_density_op} is entry \#58 of the
Calabi--Yau tables of Almkvist, van Enckevort, van Straten and
Zudilin~\cite{aesz2005} (entry 2.9 of the Calabi--Yau operator
database~\cite{cydb}), listed there as a Hadamard product, with the same
holomorphic coefficients;\ Gegelia and van
Straten~\cite{gegelia_vanstraten2024} list it among their conjectural
sources of paramodular forms, at conductors $2145$ and $2465$ (their
Table~6;\ their Table~3 lists entry~2.9 at conductors $945$ and $960$).  For every
$n$ from $2$ to $5$ the holomorphic solution of the equation of $F_n$ at
$t=\infty$ is $t^{-1}\sum_m\binom{2m}{m}^2W_{n-1}(2m)\,(16t^2)^{-m}$, the
Hadamard product of $\binom{2m}{m}^2$ with the even moments $W_{n-1}(2m)$
of the $(n-1)$-step planar random walk.  The Hadamard construction, and the
tables that place its $n=4$ and $n=5$ cases at entries \#58 and \#39, are
Almkvist and Zudilin's~\cite{almkvist_zudilin2007};\ the links of $W_N$ to
Bessel moments, random walks and Calabi--Yau equations are
Broadhurst's~\cite{broadhurst2009walks};\ Vanhove's operators for the
equal-mass two-dimensional banana integral with $N$
propagators~\cite{vanhove2014} annihilate $\sum_k W_N(2k)s^{-k-1}$, checked
here in exact arithmetic for $N=2$ to $5$, so that series is the banana's
holomorphic period at large $s$.  The
alternation of the monodromy, symplectic for $n$ even and orthogonal for
$n$ odd, matches the self-duality of the symmetric-power Kloosterman
motives of Fres\'an, Sabbah and Yu~\cite{fresansabbahyu2023}, whose
pairing is symplectic when the number of Bessel factors is even and
orthogonal when it is odd.  \textbf{[SYMBOLIC]} So the tree ladder is tied to the banana ladder,
where an $l$-loop equal-mass banana integral carries a Calabi--Yau
$(l-1)$-fold~\cite{bonisch2021}, by one exact operation:\ with
$y=1/(16t^2)$ and $F_n=t^{-1}g(y)$, the operator of $g$ is the Hadamard
product of the Legendre operator with the operator of the banana with $n-1$
propagators, checked in exact arithmetic for $n=2$ to $5$.  Since
$\int_0^\infty J_0(px)\,dp=1/x$, the same fact reads
$F_n(t)=\int_0^\infty dp\int_0^\infty x\,J_0(px)\,K_0(tx)K_0(x)^{n-1}dx$:\
$F_n$ is the $n$-propagator banana, with masses $(t,1,\dots,1)$, integrated
over a line of external momentum.  It sits on the rung of the banana with
$n+1$ propagators but is not that family:\ its singular values of $t^2$,
$\{1\},\{4\},\{1,9\},\{4,16\}$ for $n=2,\dots,5$, are the thresholds of the
banana with $n-1$ propagators.

\textbf{[SYMBOLIC + MEASURED]} At $t=1$ the ladder is the Ising class of
Bessel moments, $F_n(1)=c_{n,0}=\int_0^\infty K_0(x)^n\,dx$~\cite{bbbg2008}.
The singular values of $t^2$ are the thresholds $(n-1-2j)^2$ for every $n$
up to $8$, with no other finite nonzero singular point, so $t=1$ is a singular point
of $F_n$'s operator exactly when $n$ is even;\ the operators for $n=6$, $7$
and $8$ are irreducible by the exponent argument above, so these are the
singular points of $F_n$'s minimal operator.  There one local exponent is
doubled, while at the thresholds of odd $n$ one is a half-integer.  Since
$F_n$ is analytic for $t>0$, at a threshold it lies among the solutions
without the logarithm or the branch.  A prediction registered and committed before the
computation (\texttt{docs/predictions/2026-10-17\_ising\_class.md}) read
this as a rule for values:\ odd moments stay on their rung and even ones
drop.  The registered value cases hold, with moments for $n=2$ to $8$ to
286 digits by two routes and every informative integer-relation search
controlled;\ the registered attachment of $c_{4,0}$ to the conifold fibre
was not tested.  $c_{3,0}$ is a complex-multiplication elliptic value;\
$c_{5,0}$ and $c_{7,0}$ fit no tested lower basis, at coefficient bounds
of $10^{14}$ to $10^{17}$ and $10^{7}$ to $10^{9}$;\
$c_{4,0}=4\pi^2L(f,2)$ with $f=\eta(2\tau)^4\eta(4\tau)^4$ of weight four
and level eight, checked to $5\times10^{-261}$;\ and
$c_{6,0}=108\pi^2L(g,4)$ with $g$ of weight six and level 24, to the same
precision (both on the prediction record's computation, not a regression test).  But $c_{8,0}$ fits no single $L$-value tried, with no Tate
relation to coefficient bound $10^7$ and none with an $L$-value to
$10^{12}$, so the $n=8$ search gives parity no support as a rule for values.  \textbf{[PRIOR ART]} The $L$-value forms of $c_{3,0}$ and $c_{4,0}$ are
collected by Lim, Tu and Yu~\cite{lim_tu_yu2026}, who credit them to earlier
work, through the hypergeometric forms of Ref.~\cite{bbbg2008};\ the
numerical form of $c_{6,0}$ is theirs, and in their account of the
underlying Kloosterman motives the moments at $n=5$, $7$ and $8$ enter
only through $2\times2$ determinants of periods.  The singular points themselves are the
thresholds and pseudo-thresholds of the equal-mass banana operator
(Ref.~\cite{bonisch2021}, Eq.~(3.7));\ the threshold reading of the
values is what the ladder adds.  In the bookkeeping of Papers~18
and~34~\cite{loutey_paper18,loutey_paper34} these values lie beyond the
genus-graded embedding tier:\ $c_{3,0}$ is a rational multiple of $2^{-2/3}\Gamma(1/3)^6/\pi$, a
complex-multiplication period, and $c_{4,0}$ and, numerically, $c_{6,0}$
are $\pi^2$ times critical $L$-values of modular forms, periods of the
Kloosterman motives.  The nearest
physics analogue of a ceiling set by the calculation rather than the
system is the bound Bourjaily, McLeod, von Hippel and Wilhelm argue for on the
rigidity of massless four-dimensional Feynman integrals, $2(L-1)$ at $L$
loops~\cite{bourjaily2019}:\ there the loop order caps the geometry, here
the number of densities one integral links.  On the chemistry side,
one-center three- and four-electron correlated integrals have closed
forms~\cite{fromm_hill1987,padhy2017}, consistent with the rung needing
two-center densities.  A targeted search found no source placing a
molecular integral on these operators;\ that is an absence from a search,
not a proof of absence.  Matuzas~\cite{matuzas2026} relates one channel of
a different four-electron correlation integral, a ring of Coulomb links
with one diagonal over Slater functions, to another Calabi--Yau operator,
AESZ-330, by an inhomogeneous equation.

\subsection{A ring of three densities:\ the loop does not climb}
\label{sec:density_ring}

Every rung above is a tree.  An integral sits on a rung when the minimal
annihilator of its fibre has no composition factor above that rung.  A
ring of kernels leaves one momentum that the delta functions behind
Eq.~\eqref{eq:tree} do not remove, and one might expect a loop to raise the
class faster than a leg does.  The test case is the Coulomb triangle
$T=\langle\rho_a(1)\rho_b(2)\rho_c(3)/(r_{12}r_{13}r_{23})\rangle$, three
densities with every pair joined by a kernel.  Its explicitly correlated
relative $\langle f_{ik}f_{jk}/r_{ij}\rangle$ is the one three-body class of
the LiH correlated engine that does not separate.  The prediction, written
and committed on 2026-10-17 before the computation
(\texttt{docs/predictions/2026-10-17\_density\_loop.md}), was that the loop
does not climb.

\textbf{[SYMBOLIC + MEASURED]} The loop integrates out exactly.  In three
dimensions three massless propagators of unit index are unique, and the
uniqueness (star-triangle) relation gives
\begin{equation}
  \int\frac{d^3l}{|l-x_1|^2\,|l-x_2|^2\,|l-x_3|^2}
  = \frac{\pi^3}{|x_1-x_2|\,|x_2-x_3|\,|x_3-x_1|},
  \label{eq:star_triangle}
\end{equation}
which a Feynman-parameter evaluation confirms at five random triples to
$1.2\times10^{-12}$.  The triangle becomes an integral over the same two momenta and
three densities as the three-electron chain.  For spherical densities the
magnitude measure cancels the weight, and
\begin{equation}
  T = \frac{1}{\pi}\int_{|p-q|\le w\le p+q}
      \tilde\rho_a(p)\,\tilde\rho_b(q)\,\tilde\rho_c(w)\,dp\,dq\,dw .
  \label{eq:ring_reduced}
\end{equation}

\textbf{[SYMBOLIC + MEASURED]} At one center, with $1s$ Slater densities,
Eq.~\eqref{eq:ring_reduced} reproduces Remiddi's closed
form~\cite{remiddi1991}, as written in Ref.~\cite{ppr2004} and
differentiated in the exponents, to $2\times10^{-21}$.  Where the chain is
rational (Sec.~\ref{sec:three_electron}), the triangle carries a
dilogarithm.  At exponents $(1,1,1)$
\begin{equation}
  T = \frac{11}{3}-\frac{\pi^2}{8}-\frac32\,\mathrm{Li}_2\!\left(-\tfrac12\right)
      -\frac{46}{9}\ln 2-\frac34\ln^2 2+\frac32\ln 2\ln 3 ,
  \label{eq:ring_one_center}
\end{equation}
and at $(1,2,3)$ a 100-digit integer-relation search finds a closed form in
$\mathrm{Li}_2(1/3)$, $\mathrm{Li}_2(1/5)$, $\mathrm{Li}_2(1/6)$, $\pi^2$
and products of logarithms.  The ring raises the weight from zero to two.
In the bookkeeping of Papers~18 and~34 these one-center values are
embedding-tier constants of weight two, one step above the weight-one
two-center seed set;\ the $\pi^3$ of Eq.~\eqref{eq:star_triangle} is the
loop's normalization.

\textbf{[SYMBOLIC + MEASURED]} At the fingerprint fibre,
$\tilde\rho_v=\Delta_v^{-3}$ with scales $(1,1,t)$, the triangle is a moment
of three Struve functions,
\begin{equation}
  G(t)=\frac{\pi}{2}\,t^2\int_0^\infty x^2\,\psi(x)^2\,\psi(tx)\,dx,
  \qquad \psi=\mathbf{L}_{-1}-I_1 ,
  \label{eq:ring_struve}
\end{equation}
where $\psi$ is one density with its kernel in position space and solves
$(\theta^2-1-x^2)\psi=-2/\pi$.  The Mellin route of
Sec.~\ref{sec:four_density} gives
\begin{equation}
  L_G = -4\theta(\theta-3)^2 + t^2\,(5\theta^3-7\theta^2-5\theta+3)
        - t^4\,\theta(\theta+2)(\theta+3),
  \qquad L_G\,G = -3t^2 ,
  \label{eq:ring_op}
\end{equation}
with $\theta=t\,d/dt$, checked by quadrature of the Struve form at six
values of $t$ to 18 digits, and independently on the momentum form, with the
$t$-derivatives taken analytically.  The source term is the operator applied to a constant:\ $L_G\,1=3t^2$,
so $L_G(G+1)=0$, and the constant is what the regularized Mellin transform
cannot see.  This operator factors.
The solution of $L_G$ regular at $t=0$ is algebraic,
$(1-t^2/4)^{-3/2}$.  Write $L_G=M\,(\theta-\lambda_G)$ with $\lambda_G=3t^2/(4-t^2)$;\ the
division is exact.  The quotient $M$ is equivalent to the Legendre equation
$\theta^2-t^2(\theta+1)^2$ of $F_2$, the two-density period:\ the map
$f\mapsto t^3(1+\theta)f/(4-t^2)$ sends every solution of the Legendre
equation to a solution of $M$, an identity checked in exact rational
arithmetic.  And $G$ needs the elliptic piece:\ $(\theta-\lambda_G)(G+1)$ is a
nonzero solution of $M$, equal to $-0.161$ at $t=0.5$ and $-2.33$ at
$t=1.3$.  Neither side has a relation to the K3 family of $F_3$.
So $G$ is built from an elementary constant, an algebraic function and an elliptic
period integrated once against an algebraic kernel.  The
triangle sits on the elliptic rung of a single repulsion integral.  That is
one rung below the three-density chain.

\textbf{[SYMBOLIC + MEASURED]} The chain $C(t)$ at the same fibre, treated the same
way, is the control.  The Mellin route gives it an operator of order seven,
which divides exactly as $N_2\,R_3\,R_1\,(\theta-2)$.  The two lowest pieces
are elementary:\ $\theta-2$ has the solution $t^2$, and
$R_1=(\theta-3)-t^2(\theta+1)$ has the solution $t^3/(1-t^2)^2$.  The piece
$R_3$, of order three, is equivalent to the operator
$4\theta^3-t^2(\theta+1)^3$ of $F_3$:\ the map
\begin{equation}
  f\;\mapsto\;\frac{t^3(t^4-18t^2-4)\,f+t^3(2t^4-43t^2-40)\,\theta f
  +t(t^6-27t^4-96t^2+32)\,\theta^2 f}{(t^2-4)^2}
  \label{eq:chain_k3_map}
\end{equation}
sends every solution of $F_3$'s operator to a solution of $R_3$, an identity
checked in exact rational arithmetic.  The top piece $N_2$, of order two, is
not needed:\ $R_3R_1(\theta-2)$ already annihilates the chain, with a
residual of at most $2\times10^{-31}$ at 32 digits and $4\times10^{-19}$
at 20 digits at four values of $t$ between $0.5$ and $3.5$, while changing
one coefficient of $R_3$ gives residuals from 1.9 to 604.  The chain's only piece beyond the elementary ones is therefore
the K3 family of $F_3$.  Finer still, $R_3R_1$ splits exactly, so the order-five
operator is a left multiple of an operator $P$ of order four, and the chain
satisfies $P$:\ $PC$, divided by the solution $t^3/(1-t^2)^2$ of $R_1$, is
at most $10^{-23}$ at 20 digits and $10^{-33}$ at 30 digits at four values
of $t$, and $9\times10^{-45}$ at 40 digits at one, while a $1\%$ change in
one coefficient gives $0.02$ to $0.07$;\ the next paragraph derives this.  The second solution of
$R_1(\theta-2)$, which carries $\operatorname{artanh} t$, is absent.  So
$C=c\,t^2+t^2\int^t s^{-3}(B'f)(s)\,ds$ with $f$ a solution of $F_3$'s
operator and $B'$ a rational map.  That integral cannot be done inside the
K3 family:\ no rational map splits the K3 piece off $t^2$.  This is proved
by bounding the poles and degree of any such map from the local exponents
and searching that bound exhaustively in exact arithmetic, a search that
finds every planted split.  Because $F_3$'s operator is irreducible, the
chain is therefore not a rational combination of $F_3$ and its derivatives
plus any function without a K3 piece:\ it is a K3 function integrated once
against an elementary kernel.  So the
Coulomb ring lowers the rung at the fibre and raises the weight at one
center.  In neither limit does it climb the period ladder.

\textbf{[SYMBOLIC + MEASURED]} The same holds for the physical weights.
With the Slater densities of Sec.~\ref{sec:three_electron}, the chain's
fibre at zero separation, on the slice with scales $(t;1,1)$ and $t$ on
the central density, is
\begin{equation}
  G(t)=\int_0^\infty x\,K_1(tx)\,J(x)^2\,dx,\qquad
  J=3\int_0^x K_0-3xK_0-x^2K_1 .
  \label{eq:chain_slater_slice}
\end{equation}
The Mellin route gives an operator of order seven, which divides exactly
as $N_2\,R_3\,R_1\,(\theta+2)$.  The two lowest pieces are elementary:\
$\theta+2$ has the solution $t^{-2}$, and $R_1=(1-t^2)\theta-(1+7t^2)$ has
the solution $t/(1-t^2)^4$.  The piece $R_3$, of order three, is
equivalent to $F_3$'s operator $4\theta^3-t^2(\theta+1)^3$:\ an explicit
order-two map, rational in $t$ with poles at $t=0,\pm2$, sends every
solution of $F_3$'s operator to a solution of $N_2R_3$, and $R_3$ is the
annihilator of the image, both checked in exact rational arithmetic.
Again $N_2$ is not needed:\ $R_3R_1(\theta+2)$ annihilates $G$ with a
residual of at most $2\times10^{-31}$ at 32 digits at five values of $t$
between $0.5$ and $3.5$, while changing one coefficient of $R_3$ by one part
in 198 gives residuals from $10^3$ to $4\times10^4$.  And $G$ needs the K3
piece:\ $R_1(\theta+2)G$ is nonzero, $191$ at $t=0.5$ and $8.9$ at $t=3.5$.
A search modulo a prime finds no rational map that splits the K3 piece
off, in a box larger than the bound set by the leading terms of the local
exponents, and the same call finds a planted split.  So $G$ is an
elementary function plus a K3 function integrated once against an
elementary kernel.  The second solution of $R_1(\theta+2)$, which carries
$\operatorname{artanh}t$, is absent, as it is for the chain.  In
$w=t^2$, $\delta=w\,d/dw$, the exact system $R_1S=1$ modulo $R_3$ has a
unique rational solution $S=s_0+s_1\delta+s_2\delta^2$ inside the bound
set by the local exponents, with denominator $w(w-1)^4(2w+7)$.  So
$R_3R_1(\theta+2)=Q_1P$ with $P=(\theta+2)-S\,R_1(\theta+2)$ of order four,
and $P$ kills $t^{-2}$.  On $G$, $P$ leaves at most $2\times10^{-26}$ of its
largest term at 25 digits at $t=1/2$ and $13/10$, while a $1\%$ change in
$s_0$ gives $10^{-3}$.  Both annihilations are derived, not only measured.
With $W=R_3R_1(\theta+2)$, the Mellin operator is $N_2W$, and $N_2WG=0$ holds
with no source term.  Each term $t^{2i}Q_i(\theta)G$ of that operator is an
integral on the line $\operatorname{Re}s=c-2i$, $2<c<4$;\ moving it to $c$
crosses the poles of the Mellin transform of $G$ at $s=4,6,8$, which come from
the series of $J^2$ at $x=0$.  The six residues cancel exactly, though
none of them vanishes alone.  The recurrence itself has no boundary term, since the
relation for $J^2$ kills constants.  The same bookkeeping gives the nonzero
sources $-16/t$ for $\int_0^\infty xK_1(tx)J(x)\,dx$ and $-9216\,t^2$ for
$\int_0^\infty xK_0(tx)J(x)^2\,dx$, both equal to quadrature at 25 digits.
So $y=WG$ solves $N_2y=0$.  In the normal form $v''+qv=0$ of $N_2$,
$q=6-59/(4t^2)+O(t^{-4})$, with $q\ge5$ for $t\ge3.42$ and $q'=O(t^{-3})$.
So by Sturm comparison every nonzero real solution has infinitely many zeros
there, and its amplitude does not fall faster than $t^{-1/2}$.  A function
with an expansion in powers and logarithms of $t$ at infinity has neither
property, and $y$ has one;\ so $WG=0$.  For the chain the residues at
$s=0,2,4$ cancel the same way, and $q\to20/9$ with $q\ge2$ for $t\ge5.32$,
so $WC=0$ for its $W=R_3R_1(\theta-2)$.  Since $R_1S=1+NR_3$ with $N$ rational, $R_1P=-N\,R_3R_1(\theta+2)$, so
$P$ sends every solution of $R_3R_1(\theta+2)$ to a multiple of
$t/(1-t^2)^4$, which is odd in $t$.  Near $x=0$, $J$ is a sum of odd powers
of $x$ times powers of $\log x$.  So to the right of its strip the Mellin
transform of $G$ has poles only at $s=4,6,8,\dots$, and the expansion of $G$
at large $t$ has only even powers of $t$, times logarithms.  The
coefficients of $P$ are rational in $t^2$, so $PG$ has the same property.
On $t>1$, where $R_1$ is regular, $PG$ is a single multiple of
$t/(1-t^2)^4$, so that multiple is zero.  Since $G$ is analytic on
$\operatorname{Re}t>0$ and $P$ has rational coefficients, $PG$ is
meromorphic there;\ it vanishes on $t>1$, so $PG=0$ for all $t>0$.
Equivalently, among the solutions of $R_3R_1(\theta+2)$, $\ker P$ is the
set that are even in $t$ at infinity, and at infinity $P$ keeps exactly the
integer exponents of $R_3R_1(\theta+2)$, in $w$.  The analyticity of $G$
at $t=1$ alone, or at $t=2$ alone, does not give $PG=0$:\ $P$ has a solution
with a logarithm at $t=1$ and one with a half-integer exponent at $t=2$.
For the chain, $C=(8t^3/\pi^3)\int_0^\infty\chi(x)^2K_0(tx)\,dx$, where
$\chi=\int_0^x sK_1$ is, near $x=0$, a sum of odd powers of $x$ times powers
of $\log x$.  So $\int_0^\infty\chi^2K_0(tx)\,dx$ has only odd powers of $t$
at infinity, times logarithms, and with the factor $t^3$ the function $C$
has only even ones;\ the same steps then run with the solution
$t^3/(1-t^2)^2$ of its $R_1$.  The operator $R_3R_1(\theta+2)$ is singular at $t=1$, the equal-scale point where
the fibre collapses (Sec.~\ref{sec:three_electron});\ its exponent $-3$
there comes from $R_1$, so the singularity is not apparent.  Caps:\ the
no-split search is modulo a prime, and its bound rests on the leading-term
argument rather than the written proof the chain has;\ the annihilation
and the parity step use the standard Mellin contour shift.  Records:\
\texttt{docs/predictions/2026-10-19\_p3e2\_widened.md},
\texttt{docs/predictions/2026-10-20\_slice\_artanh.md},
\texttt{docs/predictions/2026-10-20\_slice\_parity\_proof.md} and
\texttt{docs/predictions/2026-10-21\_wg\_closure.md}.  Guard:\
\texttt{tests/test\_paper63\_chain\_slater\_slice.py}.

\textbf{[SYMBOLIC + MEASURED]} At $\gamma=1$ the correlated triangle does not drop.  Put
the exponential geminal $f=e^{-\gamma r}$ on the two edges at one density,
$Y=\langle\rho_i\rho_j\rho_k\,f_{ik}f_{jk}/r_{ij}\rangle$, at the fibre with
scales $(1,1,t)$ and $t$ on the density that carries both geminals.  The
registered prediction was that $Y$ stays on the elliptic rung;\ it fails.  The loop is
still elementary:\ it is a mass derivative of the loop with masses
$(\gamma,\gamma,0)$, a solid angle,
$\pi^2\arctan(\sqrt{\det\mathcal{G}}/\mathcal{D})/\sqrt{\det\mathcal{G}}$ in the Gram data $\mathcal{G}$, $\mathcal{D}$ of the masses
and momenta, matched by direct integration at five triples to
$3\times10^{-14}$.  But it has no uniqueness step.  The two densities and
the two geminals around it form a kite, whose Taylor series satisfies an
order-two operator, measured on 600 coefficients and mapped exactly to the
Legendre equation:\ the kite is elliptic.  At $\gamma=1$ the minimal
annihilator of $Y$ has order four and divides exactly as $\tilde R_3\tilde R_1$,
with factors of orders three and one that are not those of the chain.  The
piece $\tilde R_1$ has an algebraic solution.  The piece $\tilde R_3$ is the adjoint of an
operator whose solutions are the image of the solutions of $F_3$'s operator
under an exact order-two map, and since $F_3$'s operator is anti-self-adjoint,
$\tilde R_3$ is equivalent to it.  Applied to $Y$, the order-four operator leaves at
most $2.3\times10^{-26}$ of its largest term at four values of $t$, with no
source term, and no operator of order three annihilates $Y$:\ an exact
search within a bound set by the local exponents finds no split.  Attaching the
last density turns the elliptic kite into the K3 family, a Clausen-type
step.  So at $\gamma=1$ the correlated triangle sits with the three-density chain, on the
K3 rung, and not with the Coulomb triangle.  At $\gamma=3/2$ an annihilator of $Y$ of order nine
has an exact composition series whose pieces have orders one (four of
them, each with an algebraic solution), two and three.  The order-three
piece has the local data of $F_3$'s operator twisted by $t$, with genuine
singular points only at $t=0,\pm2,\infty$, which fixes it up to
equivalence by Levelt's rigidity, in the form of Beukers and
Heckman~\cite{beukers_heckman1989};\ an exact order-two map from $F_3$'s
operator confirms the equivalence.  The order-two piece is elliptic.  In
$x=t^2$ it is equivalent, by an exact map of order one, to the operator of a
period of the rational elliptic surface
$Y^2=X^3-3x^2(16x^2+152x+25)\,X-2x^3(4x-5)(16x^2-616x+25)$, whose
$j$-invariant is $(16x^2+152x+25)^3/\bigl(x(4x-1)^2(4x-25)^2\bigr)$ and whose
singular fibres are $I_1^*$, $I_2$, $I_2$, $I_1$ at $x=0,1/4,25/4,\infty$;\ over
$t$ they are six fibres $I_2$, at $t=0,\pm1/2,\pm5/2,\infty$, that is at
$t=|1-\gamma|$ and $1+\gamma$.  The monodromy of the piece has integer traces,
measured at 80 digits, and the map is exact (record
\texttt{docs/predictions/2026-10-22\_c2b\_identification.md}).
\textbf{[SYMBOLIC + MEASURED]} The curve is Legendre's, with modulus
$\lambda=(\gamma^2-(t-1)^2)/(4t)$, so that $\lambda(1-\lambda)$ is
$-1/(16t^2)$ times the K\"all\'en function of $t^2$, $1$ and $\gamma^2$.
Twisted by $t$, the Legendre operator of degree $-1/2$ in
$z=(t^2+1-\gamma^2)/(2t)$ is the period operator above, coefficient for
coefficient.  \textbf{[PRIOR ART]} The function itself is
$\pi\bar S_3(a,b,w)$ of Ref.~\cite{bbbg2008} (their Eq.~30), given there
as a convolution of two square-root propagators,
$(\pi/4)\int_{-\infty}^{\infty}dk/\sqrt{(k^2+a^2)((k-c)^2+b^2)}$, and as a
complete elliptic integral of modulus
$\sqrt{(c^2+(a-b)^2)/(c^2+(a+b)^2)}$;\ the Legendre form, the imaginary
momentum and the triangle are the readings added here.  It annihilates
$C(t)=\int_0^\infty K_0(tx)K_0(x)\cosh(\gamma x)\,dx$, which converges for
$t>\gamma-1$, the pseudo-threshold, to $5\times10^{-41}$ at 40 digits, and $C=\pi^2P_{-1/2}(z)/(4\sqrt t)$ to $2\times10^{-16}$ at
15 digits:\ the continuation to imaginary momentum of a classical cosine
transform of $K_0K_0$.  Since $\int dp_2\,e^{ip_2x_2}=2\pi\delta(x_2)$,
$C$ is $\pi$ times the two-dimensional bubble with masses $t$ and $1$,
integrated over a line of external momentum moved off the origin by
$i\gamma$;\ at $\gamma=0$ it is $F_2$.  In the dictionary of
Sec.~\ref{sec:four_density} the geminal therefore acts as an imaginary
external momentum, and the singular points are where $\gamma=t+1$ or
$\gamma=|t-1|$, the bubble's threshold and pseudo-threshold.  The
registered candidate with $e^{-\gamma x}$ in place of $\cosh(\gamma x)$
is not a solution, because its sinh part is not (record
\texttt{docs/predictions/2026-10-22\_c2b\_feynman.md}).  This fixes the
operator class of the piece, not the component of $Y$ along it.
\textbf{[SYMBOLIC + MEASURED]} At $\gamma=2$ the reading was registered
before $Y$ was computed there, and it held piece by piece.  An
annihilator of order nine, fitted from 800 series coefficients and
checked on $Y$ to $5\times10^{-26}$, has the same composition series
as at $\gamma=3/2$.  Its order-two piece is equivalent by an exact map
to the Legendre operator with $z=(t^2-3)/(2t)$, its order-three piece is
the twisted K3 piece at $t=\pm2$ again, and its order-one pieces are
algebraic, branched at $t=2$ and $t=2\gamma=4$ (record
\texttt{docs/predictions/2026-10-22\_c2b\_feynman\_g2.md}).
\textbf{[SYMBOLIC + MEASURED]} Why the order drops at $\gamma=1$.  In
Feynman parameters the kite's density pair enters through
$M(u)=\sqrt{1+u(1-u)z}$ and its geminal pair through
$N(v)=\sqrt{\gamma^2+v(1-v)z}$.  At $\gamma=1$ the two are the same
function:\ the two bubbles have equal masses, and the four pair thresholds
$t=2$, $2\gamma$, $1+\gamma$, $|1-\gamma|$ fall onto $2$ and $0$.  Two
readings were registered and separated before computing:\ equal masses,
or merely colliding thresholds.  At $\gamma=3$ the pseudo-threshold
$|1-\gamma|$ falls on the density-pair threshold $t=2$ with unequal
masses.  There the annihilator of $Y$ keeps order nine and the series
$1,1,1,1,2,3$, the elliptic piece is present, mapped exactly from the
Legendre operator with $z=(t^2-8)/(2t)$, and the genuine singular points
are exactly $t=2,4,6$;\ at $t=2$ the K3 piece's half-integer exponents and
the elliptic piece's $I_2$ fibre sit together, and no factor is lost.  At
$\gamma=1/2$, before the branch point crosses the contour, the order is
nine as well, with the Legendre map at $z=(t^2+3/4)/(2t)$ and singular
points at $t=1/2,1,3/2,2$.  So colliding thresholds remove nothing, and
the drop at $\gamma=1$ is the equal-mass degeneration, the step that
makes the equal-mass sunrise special in the Feynman-integral
literature~\cite{blochvanhove2015}.  The kite's operator in $z$ has order
six at $\gamma=1/2$, $3/2$, $2$ and $3$, with genuine singular points
exactly at minus the squared pair masses and at $0$ and $\infty$, and
order two at $\gamma=1$.  The annihilators are checked on $Y$ to
$6\times10^{-27}$ at $\gamma=1/2$ and $3\times10^{-24}$ at $\gamma=3$,
and the maps are exact (record
\texttt{docs/predictions/2026-10-23\_gamma1\_mechanism.md};\ guard
\texttt{tests/test\_paper63\_mg\_mechanism.py}).  Cap:\ the
mechanism is read from the integrand and tested by its consequences at
two values of $\gamma$;\ no symbolic derivation of the drop from $N=M$ is
given.
\textbf{[SYMBOLIC + MEASURED]} The cusp-correct geminal $f=re^{-\gamma r}$
of the LiH engine, on both edges at $\gamma=3/2$ and $2$, brings no new type, as
registered before computing from the fact that $re^{-\gamma r}$ is a
$\gamma$-derivative of the same Yukawa triangle.  Its annihilator of order
nine, checked on the integral to $10^{-28}$, has the same composition
series:\ four algebraic pieces, branched at $t=2$ and $3$;\ an order-two
piece, a different operator from the exponential case's, equivalent by an
exact map to the same Legendre family;\ and the twisted K3 piece (record
\texttt{docs/predictions/2026-10-22\_geminal\_shapes.md};\ guard
\texttt{tests/test\_paper63\_ga1\_geminal.py}).  At $\gamma=2$,
registered as a check of the rule, the order is nine again with the same
series and the same degrees, the genuine singular points are exactly
$t=1,2,3,4$, the order-two piece is again a different operator from the
exponential case's and maps exactly from the Legendre operator with
$z=(t^2-3)/(2t)$, and the algebraic pieces carry the exponents $-11/2$ and
$-13/2$ at $t=2\gamma$ as at $\gamma=3/2$;\ the operator is exact, from
coefficients modulo twelve primes and rational reconstruction, and leaves
$2\times10^{-24}$ on the integral at 32 digits (record
\texttt{docs/predictions/2026-10-23\_geminal\_shapes\_g2.md};\ guard
\texttt{tests/test\_paper63\_ga3\_geminal.py}).
\textbf{[SYMBOLIC + MEASURED]} One function carries both elliptic pieces
met so far.  Write
$\mathcal{L}(a,b;c)=\int_0^\infty K_0(ax)K_0(bx)\cos(cx)\,dx
=\pi^2P_{-1/2}(z)/(4\sqrt{ab})$, with $z=(a^2+b^2+c^2)/(2ab)$.  At $c=0$ it
is the radial core of the three-center integral at zero separation:\
with $a_i=c_i^{-1/2}$ it reproduces Eq.~\eqref{eq:period} to
$4\times10^{-31}$ at three pairs of scales.  At $c=i\gamma$ it is $C(t)$
above.  Zero or real momentum gives $z\ge1$.  Imaginary momentum gives
$z=\cos\theta$ between the pseudo-threshold and the threshold,
$|t-1|<\gamma<t+1$, with $\theta$ the angle opposite the side $\gamma$ in
the triangle with sides $t$, $1$ and $\gamma$, and the modulus is
$\lambda=\sin^2(\theta/2)$.  The thresholds are the flat triangles, since
the K\"all\'en function is $-16$ times the squared area.  The slot is not
a convention:\ with the geminal as a mass and a density scale as the
momentum, the two other assignments give the families
$j=4(16x^2+472x+25)^3/\bigl(9x(4x-1)^2(4x-25)^2\bigr)$ and
$j=4(16x^2-104x+601)^3/\bigl(9(4x-1)^2(4x-25)^2\bigr)$, and neither is
isogenous to the piece's, since their traces of Frobenius differ at most
good primes at every rational $x$ tried, while the two solutions of the
degree-six search agree at every one.  The right angle,
$\gamma^2=t^2+1$, gives $\lambda=1/2$ and $j=1728$, the same
complex-multiplication value $K(1/\sqrt2)$ as the three-center fibre at
$\rho=1/2$.  \textbf{[OBSERVATION]} On this family, then, composing
unequal decay rates costs one elliptic period:\ at zero momentum its
modulus is set by their ratio, and equal rates make it elementary, as
Eq.~\eqref{eq:period} records;\ at imaginary momentum it is set by the
angle of their triangle, and a flat triangle is singular.  Reading
one-loop integrals through triangles of masses and momenta is the
approach of Davydychev and Delbourgo~\cite{davydychev_delbourgo1998};\
what is added is its appearance in a correlated molecular integral.
\textbf{[SYMBOLIC + MEASURED]} The shape of the geminal matters at this
level.  A Gaussian geminal's avatar, $\int_0^\infty K_0(tx)K_0(x)e^{-x^2}dx$,
is annihilated by $t^4(\theta+1)^2-t^2\theta^2(4\theta+1)+4\theta^2(\theta-2)^2$
to $10^{-41}$, and its monodromy companion, with $I_0$ for $K_0(tx)$, grows
like $e^{(t-1)^2/4}$, so every annihilator has an irregular singular point
at infinity:\ at the level of the one function, Slater geminals stay on the
ladder of Picard--Fuchs operators and Gaussian ones leave it for exponential
periods (registered before computing, record
\texttt{docs/predictions/2026-10-22\_geminal\_shapes.md}).
\textbf{[SYMBOLIC + MEASURED]} On the full triangle, with $e^{-r^2}$ on
both edges and scales $(1,1,t)$, the Gaussian shape removes more than the
elliptic piece.  The minimal annihilator of $Y$ has order nine and
$x$-degree 13 and splits exactly into two pieces, of orders four and five,
both irregular at infinity, with formal solutions growing like
$e^{t^2/4}$ and $e^{t^2/8}$;\ its only genuine finite singular point is
$t=2$.  No piece is equivalent to $F_3$'s operator, although each keeps
the K3 piece's Jordan block of size three at $t=0$:\ the K3 rung of the
Slater triangle is not set by the densities and the Coulomb line alone.
The prediction that the K3 piece survives failed.  The operator leaves
$2\times10^{-27}$ on $Y$ at 32 digits, and an independent evaluation of
$Y$ agrees to $4\times10^{-14}$ (record
\texttt{docs/predictions/2026-10-23\_gaussian\_geminal\_triangle.md};\ guard
\texttt{tests/test\_paper63\_gy\_gaussian.py}).
Caps:\ one value, $\gamma=1$, the equal-mass value, at which the Slater
triangle has no elliptic piece either, so the comparison of shapes is
cleanest for the K3 piece;\ the irreducibility of the two pieces and the
absence of an $F_3$ piece rest on bounded searches;\ and minimality assumes
that $\sqrt{\pi/2}\,e^{1/2}\operatorname{erfc}(1/\sqrt2)$, the constant
that enters the Taylor coefficients, is irrational.
Guards:\ \texttt{tests/test\_paper63\_one\_formula.py},
\texttt{tests/test\_paper63\_geminal\_shapes.py}.  The order-six right factor below the order-three piece
leaves between 1.3 and 2.4 times the largest term of $Y$ at three values
of $t$, so $Y$ needs the K3 piece, and no piece above K3 occurs (record
\texttt{docs/predictions/2026-10-22\_correlated\_triangle\_g32.md}).  The singular points of $Y$'s annihilator move with
$\gamma$, to $t=|1-\gamma|,2,1+\gamma,2\gamma$, three of which meet at $t=2$ for
$\gamma=1$.  \textbf{[OBSERVATION]} The drop of the Coulomb ring is therefore read as a property of
uniqueness, not of rings:\ both loops integrate out exactly, but only the
Coulomb one factors into pairwise inverse distances, and only it lowers the rung.

\textbf{[OBSERVATION]} The registered decision rule did not anticipate this
outcome.  It allowed the Coulomb triangle to land on the chain's rung, with or
without an elementary term, or to climb;\ it lands below.  By the rule's
letter $L_G$ is an order-three factor that is not equivalent to the chain's
operator, which the rule called a climb, but $L_G$ factors into lower pieces,
so it is not a climb in the sense the record intended.  Two side readings
failed.  The singular points at $t=\pm1$ were not predicted;\ they belong to
the elliptic piece.  And a Coulomb ring of three does not sit with the chain of
three, as the record expected, but lower.  The count of linked densities
orders the trees, not the rings.  The caps of Sec.~\ref{sec:four_density}
apply:\ these are fibre statements, and the integrated value may sit lower.
Further caps:\ for the correlated triangle, the exponential shape at
five values of $\gamma$, the first the equal-mass value, with $t$ on the
density carrying both geminals, and the LiH engine's cusp-correct shape at
$\gamma=3/2$ and $2$ only, and a Gaussian shape at $\gamma=1$ only;\ the kite operators were guessed from series and
checked on their coefficients, and are now also derived by integration by
parts (below), and the mechanism of the drop at $\gamma=1$
is tested by its consequences at two values of $\gamma$, not derived;\ only
the triangle among rings;\ and that the chain satisfies its own order-four $P$ rests on the
Mellin contour shift.  The search that excludes a K3 piece
from the triangle is bounded in degree, unlike the chain's no-split proof,
but the orders settle it:\ the
factors found have orders one and two, which leaves no room for the
order-three K3 factor.
Guards:\ \texttt{tests/test\_paper63\_density\_ring.py}.

\textbf{[PRIOR ART]} The one-center values are Remiddi's~\cite{remiddi1991}:
the three-electron correlation integral over Slater functions closes in
logarithms and dilogarithms, and Fromm and Hill~\cite{fromm_hill1987} treat
the general case.  The uniqueness relation~\eqref{eq:star_triangle} is
classical;\ its general-dimension form, which gives the constant $\pi^3$ at
$D=3$, and the early references are collected by Preti~\cite{preti2018}.
Under the relation of Sec.~\ref{sec:four_density} to the banana ladder, the
triangle's elliptic piece is the $n=2$ rung itself.  What is added here is the fibre reading:\ the Struve
moment~\eqref{eq:ring_struve}, its operator, its place below the chain,
and the correlated triangle's place with it.  The solid-angle form of the massive loop is Davydychev and Delbourgo's
(Sec.~5.1 of~\cite{davydychev_delbourgo1998}, which recovers Nickel's
result~\cite{nickel1978});\ here it also rests on its numerical check.
No literature search was made for the operator~\eqref{eq:ring_op}.

\textbf{[SYMBOLIC + MEASURED]} The kite operators above were guessed from
Taylor series.  They also follow from integration by parts, the standard
reduction of Feynman integrals~\cite{chetyrkin_tkachov1981,laporta2000}.  In
momentum space the kite is a two-loop integral in $D=3$ with five lines:\ the
two densities with index $3/2$, the two geminals $(q^2+\gamma^2)^{-2}$ and
the Coulomb line $q^{-2}$.  A reduction modulo a large prime, sampled at many
values of $z$, gives the differential system of the master integrals, and
the kite's cyclic vector in that system gives its operator.  At
$\gamma=1/2$, $3/2$, $2$ and $3$ the operator has order six and equals the
guessed one exactly.  At $\gamma=1$ the master integrals drop from nine to
three and the operator is the Legendre one, so the equal-mass degeneration
shows in the reduction itself.  On the tetrahedron, the three-loop vacuum
integral with the apex density of mass $t$, the reduction gives an operator
for $Y$ of order ten at $\gamma=3/2$, divisible on the right by the
order-nine operator above.  Off the symmetric slice, at density scales
$(1,2)$ and geminal exponents $(3/2,3)$, the series route found no operator
with 260 unknowns or fewer.  The reduction gives one of order eleven and
degree $113$ in $z$, about 1300 unknowns.  Its integer coefficients reach 565
bits, but the eighteen-dimensional system behind it needs only 37, so the
system is rebuilt exactly from two primes and checked on a third, and the
operator is built from it over $\mathbb{Q}(z)$.  It agrees with the operator
computed at thirteen primes, and on 300 Taylor coefficients at 280 digits it
leaves $1.3\times10^{-280}$ of the largest term;\ with one coefficient
changed by $1/200$ it leaves $4\times10^{-41}$.  Its genuine singular points
are exactly $z=-1/4,-1,-9/4,-4,-9,-49/4,-16,-81/4$, with $0$ and $\infty$:\
minus the squares of the two-line cuts through the two densities and through
the two geminals, and of the three-line cuts that pair each density with the
geminal on the opposite edge.  The points $z=-25/4$ and $-25$, where a
pairing of each density with its adjacent geminal would put its cuts, are
ordinary, and the remaining 92 roots of the leading coefficient are
apparent.  Controls:\ the one-loop bubble in $D=3$ against its arctan closed
form, and the kite with the Coulomb line removed, whose exact Taylor series
is the product of two bubbles (records
\texttt{docs/predictions/2026-10-24\_ibp\_kite.md} and
\texttt{docs/predictions/2026-10-23\_price\_formula\_slots.md}).
\textbf{[SYMBOLIC + MEASURED]} On the tetrahedron at the same point the
reduction needs 34 masters and gives an operator for $Y/t^3$ in $x=t^2$ of
order nineteen, exact over $\mathbb{Q}$ (the system rebuilt from two primes
and checked on a third).  On $Y$ from a 60-digit grid it leaves
$5\times10^{-55}$ to $3\times10^{-43}$ of the largest term, following the
quadrature error.  The genuine singular points of $Y$'s minimal operator in
$x>0$ are the same eight cuts, $x=1/4,1,9/4,4,9,49/4,16,81/4$.  The other
roots of the leading coefficient are apparent, and at each of the eight some
analytic continuation of $Y$ changes on a small loop.  The change is the
same to four digits in 64- and 90-digit arithmetic, and at least $10^{10}$
times the change on a loop around an ordinary point along a path at least as
long.  \textbf{[SYMBOLIC + MEASURED]} The order nineteen is not minimal.  The
reduction treats the two sectors with one geminal and no Coulomb line as
different, but each is the same sunrise of the three densities times a
one-line tadpole, so their masters differ by the constant ratio of the
geminal scales, $\gamma_1/\gamma_2=1/2$.  With that relation imposed the
operator has order fifteen, exact over $\mathbb{Q}$, and on $Y$ it leaves
$2\times10^{-56}$ to $4\times10^{-46}$ of the largest term;\ the same
construction with the ratio $1/3$ leaves $10^{-4}$ to $2\times10^{-2}$.  The
reduction's system is block triangular in its sectors, so the pieces can be
read sector by sector, and they are the predicted ones.  The K3 operator of
the chain of three with scales $(t;1,2)$ is the sunrise sector's own block, and
it is a left factor.  The Legendre family pulled back to the two crossed
pairings, $z_1=(4t^2+7)/(16t)$ and $z_2=(t^2-8)/(2t)$, gives the two pieces of
order two, from the sectors with one geminal and the Coulomb line;\ the direct
pairings do not occur.  The other seven pieces have order one and algebraic
solutions, branched at $x=1,9$, at $9/4$ and at $81/4$, and come from the
sectors with both geminals.  Each of the three identifications is an exact map
over $\mathbb{Q}(x)$, rebuilt from many primes and checked by exact division.
$Y$ needs the K3 piece:\ the right factor without it leaves $2\times10^{-4}$ to
$4\times10^{-3}$ of the largest term.  It needs both elliptic pieces, since its
minimal operator is genuinely singular at $x=1/4$ and $x=4$ and no other piece
is singular there.  \textbf{[SYMBOLIC + MEASURED]} The minimal operator of
$Y$ has order thirteen, exact over $\mathbb{Q}$.  Two of the three order-one
pieces branched at $x=9/4$ are left factors of the order-fifteen operator.
The pairing of $Y$ with each of them, the bilinear concomitant, is a constant,
and on $Y$ it is below $10^{-56}$ of its largest term and grows with the
quadrature error;\ the same pairing with the K3 piece, the control, is
$10^{-8}$ of its largest term and does not move at twice the step.  So the
right factor of order thirteen without those two pieces annihilates $Y$,
leaving $8\times10^{-57}$ to $2\times10^{-46}$ of the largest term.  Its only
irreducible left factor is the K3 piece, once, and $Y$ needs that piece, so no
smaller operator annihilates $Y$.  Its five order-one pieces are branched at
$x=1,9$ (one), at $9/4$ (one) and at $81/4$ (three).  The K3 piece is
irreducible:\ its exponents leave no room for a factor of order one or three,
and its exterior square has no element of the one admissible form, so it has
no factor of order two.  Its monodromy preserves a symmetric form and no
alternating one, with reflections at $x=1$ and $x=9$, and its Jordan blocks
are $2+2$ at $x=0$ and $3+1$ at infinity.  Caps:\ the reduction is new code,
checked against the guessed operators and the series;\ the two vanishing
pairings are measured, not proved;\ the left-factor searches behind the order
thirteen are bounded, with planted controls found;\ and the form, the blocks
and the irreducibility are computed modulo two primes, the last by searches
that are complete for its exponents.
An earlier reading of the order-nineteen operator, as one piece of order twelve
with none of these pieces at either end, came from a map search whose box was
too small at the left end, and is withdrawn (records
\texttt{docs/predictions/2026-10-23\_price\_formula\_slots.md}, addendum 6,
\texttt{docs/predictions/2026-10-25\_geometric\_route\_y.md} and
\texttt{docs/predictions/2026-10-25\_tetrahedron\_route\_y.md}).
\textbf{[PRIOR ART]} Integration by parts on
this tetrahedron in momentum space is Puchalski, K\c{e}dziera and
Pachucki's, for the three-electron Slater integral with integer
indices~\cite{pkp2009}, which they reduce to a differential equation.  What
is added here is the family with half-integer density lines and the
geminal's squared lines, and the check against operators found
independently.  Guards:\ \texttt{tests/test\_paper63\_ib\_kite.py} and
\texttt{tests/test\_paper63\_y\_minimal\_order.py}.

\textbf{[MEASURED]} The operators are also evaluators.  The kite's operator,
with the kite's first six Taylor coefficients at $z=0$ as its only input,
gives the kite at all $247$ nodes of the grid behind $Y$.  It does so by Taylor
steps along a path that passes above the three apparent singular points on
the positive axis.  Each step reaches half the distance to the nearest root
of the leading coefficient, apparent roots included, because rounding
excites solutions of the recurrence that blow up there.  On one process this
takes $64$\,s, of which the input takes $1$\,s;\ the quadrature it replaces
took $7894$\,s on fifteen.  $Y$ from the two routes agrees to
$2\times10^{-69}$ at four values of $t$, the precision at which the grid is
stored.  The node error grows towards large $w$, to $3\times10^{-8}$ near
$w=10^{12}$, but those nodes carry so little weight that no node moves $Y$ by
more than $1.4\times10^{-68}$.  With the derivation included, the route
costs about $9300$\,s of processor time, about a thirteenth of the upper
bound on the quadrature's.  The derivation is three reduction runs on six
worker processes, one of which was measured at $3065$\,s of processor time,
and a $13$\,s lift.  Changing
one operator coefficient by $1/200$ moves $Y$ by $1$ to $7\%$.  $Y$'s own
operator does not pay in the same way.  From the grid's data at $t=1.7$ it
reaches five other values of $t$ to $4\times10^{-48}$ or better, but at $4$ to
$12$\,s per value, against $0.02$\,s for one more grid sum, because once the
kite is known the outer integral is a cheap sum.  At $60$ digits the structure pays where
the cost is, in the inner kite.  Caps:\ one parameter set, so a new
configuration needs a new derivation, which dominates the cost;\ times on one
machine, with the quadrature's processor time an upper bound;\ the
comparison is to this paper's own grid, not to a production Slater-integral
code such as Ref.~\cite{chen2026fourier};\ and no
certified error bounds, only agreement with the grid (record
\texttt{docs/predictions/2026-10-26\_operator\_evaluator.md}, registered before
computing;\ guard \texttt{tests/test\_paper63\_operator\_evaluator.py}).
\textbf{[MEASURED]} At chemistry precision the comparison reverses.  For $Y$ to $10^{-12}$, a
double-precision tanh-sinh quadrature of the same two-dimensional integral, graded at the
ends and split along its ridge, takes $0.033$\,s of processor time on one core.  The
operator route takes $16$\,s with the operator already given, about $500$ times slower,
and $3\times10^{5}$ times slower with its derivation counted.  It also cannot run below
$60$ digits:\ at $20$ to $40$ digits its Taylor walk diverges.  The two routes agree on $Y$
to $7\times10^{-14}$.  So the operator is an evaluator only where high precision is wanted;\
for chemistry it is a statement about structure, not a shortcut (record
\texttt{docs/predictions/2026-10-28\_kite\_speed\_test.md}, registered before timing;\ one
parameter set, one machine).  \textbf{[PRIOR ART]} Evaluating Feynman integrals by solving
their differential equations in series is standard in particle
physics~\cite{moriello2020,hidding2021,liu_ma2023}, and
Ref.~\cite{pkp2009} reduces the integer-index three-electron integral to such
an equation.  What is added is the measured cost on this family.

\textbf{[MEASURED]} Most of such an operator's singular points can be read
off before it is derived, but not all of them.  At a second set of rates,
chosen blind, $(a,b,\gamma_1,\gamma_2)=(1,3/2,7/3,8/3)$, the kite's operator
again has eighteen masters and order eleven.  The eight singular points
predicted from its cuts, at minus the squared sums and differences of the
rates on each cut, are all genuine and carry the predicted local exponents;
the pairings of adjacent lines give none.  But a ninth genuine point appears
at $z=-111/1120$, with exponents $-7/2,0,1,\dots,9$, and the leading Landau
equations of the kite and of its contractions have no solution there.  A
count of critical points of the Lee--Pomeransky likelihood, which equals the
number of masters at a generic point, confirms it:\ $22$ at generic $z$, $21$
there, at two primes.  The drop is in the full kite only.  Its position has a closed form in the squared rates
$A=a^2$, $B=b^2$, $G_1=\gamma_1^2$, $G_2=\gamma_2^2$:
\begin{equation}
z_9=\frac{(AG_2-BG_1)(A-B-G_1+G_2)}{(A-G_1)(B-G_2)} .
\label{eq:kite_ninth_point}
\end{equation}
It is symmetric under the kite's reflection $(a,\gamma_1)\leftrightarrow
(b,\gamma_2)$ and under $(a,b)\leftrightarrow(\gamma_1,\gamma_2)$.  Its poles
lie where an adjacent pair has equal rates, the pairings the cuts ignore.  At
the original rates $a\gamma_2=b\gamma_1$, so $z_9=0$:\ the ninth point sits on
$z=0$ there, which accounts for that operator's exponent at $z=0$ that the
cuts did not explain.  The formula was fitted from $41$ sets of rates, $29$
unknowns, and passed $18$ held-out sets at two primes;\ it reproduces
$-111/1120$ at the blind point, and on the family $(1,2,3/2,\gamma_2)$ the
critical-point count drops from $22$ to $21$ exactly at $z_9$ (record
\texttt{docs/predictions/2026-10-26\_kite\_ninth\_point.md}, registered before
computing).  \textbf{[PRIOR ART]} The location is known.  Broadhurst gives
it, with $s=-z$, as the pole $s_3$ that the two-particle cuts of the kite
acquire when its rung becomes massless with unequal adjacent masses
(Ref.~\cite{broadhurst2022kites}, Eqs.~(24)--(25));\ it is the massless-rung
limit of the anomalous threshold of the triangle formed by the two lines of
one cut and the rung, the other two lines on shell, which the kite's own
leading Landau equations do not contain.  \textbf{[MEASURED]} A registered
test confirms that this is the mechanism here.  With all five lines on shell
and the loop momenta in one plane, eliminating the equations in components
gives exactly the factor of Eq.~\eqref{eq:kite_ninth_point} at three sets of
rates and no other.  The multipliers are all nonzero, and the first Symanzik
polynomial $U$ vanishes on them, which is why a Landau test that requires
$U\neq0$ misses the point.  Given a small Coulomb mass, the point does not
vanish but moves:\ the count drops on a curve through $(z_9,0)$, one of
Broadhurst's two triangle thresholds, while the other runs to infinity as the
mass goes to zero.  A scan in $z$ that assumes no location finds, at the blind
rates, exactly ten drop points:\ $z=0$, $z_9$ and the eight cut values.  Caps:\
the massive curve is checked modulo two primes at three masses, not over the
rationals, and the local shape, the drop size on the curve and the escape
direction were not run (record
\texttt{docs/predictions/2026-10-27\_kite\_ninth\_point\_origin.md}, registered
before computing;\ guard \texttt{tests/test\_paper63\_kite\_ninth\_point.py}, which also
covers the blind point and Eq.~\eqref{eq:kite_ninth_point}).  What is added here is the measurement in this family:\ the
critical-point drop, the exponent, and the complete drop locus at one set of
rates.  The operator's degree,
$120$, exceeds the prediction, and its coefficients needed a fourth prime.
With $\gamma_2^2$ left free the differential system has degree $16$ in it,
but the kite written in the masters has degree $26$, so the one-parameter
family exceeded its registered sample cap and was not built (record
\texttt{docs/predictions/2026-10-26\_parametric\_kite\_operator.md},
registered before computing).  \textbf{[PRIOR ART]} Reading singular points
from cuts is standard Landau and threshold analysis.

\textbf{[MEASURED]} In the molecular calculation the triangle class is a
modest share of the cost.  In the explicitly correlated LiH calculation of
Paper~19 at its best truncation, the class that $Y$ models takes $21\%$ of a
one-process run, third of its four large pieces, behind the electron--nucleus
and kinetic blocks.  The part an operator evaluator would replace, the
factorization of the three-body kernels, is $3\%$ of the run;\ the rest of
the class is about $2.4$ million Coulomb mode solves.  So a free triangle
class would speed the run by at most $1.27$ times, and a free kernel build by
$1.03$ times.  Caps:\ one run per truncation on one machine;\ a reduced
truncation is no proxy, giving $48\%$ (record
\texttt{docs/predictions/2026-10-26\_lih\_triangle\_share.md}, registered
before computing).

\textbf{[MEASURED]} Whether the one-mass slice could be tabulated, as Gaussian
codes tabulate the Boys function, depends on how many master integrals it
needs and in how many variables.  Counting critical points, as for the kite,
the slice needs four masters, and the index range the production code uses
needs a fifth, an antiderivative in $D$ from a pole at $x=0$;\ a scaling
identity leaves two variables.  By the registered gate this is borderline:\ a
table is not ruled out, but its speed against a production
code~\cite{chen2026fourier} is untested.  The physical fibre, with two
offsets, needs six masters in three variables.  Caps:\ the scaling identity
missed its registered precision, $2.7\times10^{-25}$ against $10^{-25}$, and
an unregistered rerun agrees to $1.6\times10^{-37}$;\ these are counts of
masters, not the cost of a table (record
\texttt{docs/predictions/2026-10-27\_master\_count\_tabulation.md}, registered
before computing;\ guard \texttt{tests/test\_paper63\_master\_count.py}).

\textbf{[MEASURED]} Production explicitly correlated codes do not use a
single Gaussian geminal.  For many integrals they replace $e^{-\gamma r}$ by
a least-squares fit, six Gaussians by default.  The off-slice triangle above,
known to about fifty digits, measures that replacement.  With the fit in the
plain norm $\int_0^\infty r^2(\cdot)^2\,dr$, which at three and six terms gives
the STO-3G and STO-6G exponents~\cite{hehre1969}, the six-term fit makes $Y$ too small by
$3.6\times10^{-5}$ at $t=1/2$ down to $2.4\times10^{-6}$ at $t=4$.  The error
falls like $e^{-b\sqrt{n}}$ in the number of terms $n$, the law the cusp at
$r=0$ sets for Gaussian expansions of exponentials~\cite{kutzelnigg1994}, proved
as an upper bound in a weighted norm by Braess~\cite{braess1995};\ the constant
$b$ on this integral is our measurement.  With an extra weight $e^{-\gamma^2r^2}$ the fit loses its tail, and
the six-term error is $5.4\times10^{-2}$ to $1.7\times10^{-2}$.  So on this
integral the error is set by how far out the fit is trusted, not by the cusp.
Caps:\ the production weight function is not reproduced, so these two weights
do not price a production fit;\ one configuration, with model densities (record
\texttt{docs/predictions/2026-10-25\_gaussian\_fit\_geminal.md}, registered
before computing).
\textbf{[MEASURED]} A fit that a production code ships, the six-term fit
written into Psi4's source~\cite{psi4_2020}, is the best of the three at the
cusp and the worst in the tail.  It is $1.8\%$ low at $r=0$, against $6\%$ for
the plain fit, and half the true value at $\gamma r=4$.  On the same $Y$ it
gives an error of $-13\%$ at $t=1/2$ down to $-5.5\%$ at $t=4$, more than three
thousand times the plain fit's.  The three fits are ordered by their tails,
not by their cusps.  Caps:\ $Y$ weights the geminal over its whole range,
while a correlation energy is dominated by short range, so this prices the
fit on this integral and not on an F12 energy;\ the fit's weight and its
literature source are not verified, nor which integrals the code uses it for
(record \texttt{docs/predictions/2026-10-26\_production\_geminal\_fit.md},
registered before computing;\ guard \texttt{tests/test\_paper63\_production\_geminal\_fit.py}).


%% ========================================================================
\section{The graph reading, as a bounded conjecture}
\label{sec:conjecture}
%% ========================================================================

\textbf{[OBSERVATION + PRIOR ART]} Read as a Feynman graph, the cases of Paper~59's geminal section and of
Sec.~\ref{sec:density_ring} fit one pattern.  The densities and the kernels are the lines;\ a
Slater geminal is a line whose mass is its exponent $\gamma$, and the Coulomb kernel is a line
with no mass.  On every instance tested the class is set by the shape of the graph together
with the masses on its lines.  A massless Coulomb line adds no threshold, so for Coulomb
integrals the density topology alone sets the rung, as stated above.  A Slater geminal keeps
the curve the densities build, but it lifts the same-pair master above to genus one, and the
correlated triangle's operator drops order where the geminal's mass equals that of a density
pair, $\gamma=1$.  A Gaussian geminal is not a propagator of this kind, and its integrals leave
the Picard--Fuchs ladder (Sec.~\ref{sec:density_ring}).  That the internal masses set the
geometry of a Feynman integral is standard:\ the sunset graph with massless lines evaluates in
Gamma functions, and with massive lines it is elliptic~\cite{blochvanhove2015,vanhove2014}.
What is recorded here is only that these molecular integrals follow the same dictionary.  The
nearest theorem in shape is Brown's:\ for massless primitive $\phi^{4}$ graphs in four
dimensions the set of periods is minor monotone, so a minor-closed class of graphs bounds the
ring the periods lie in~\cite{brown2009periods}.  That theorem has no masses;\ here the masses
enter too.  It is
a pattern across the tested instances, not a derivation.  It would fail on a Coulomb-only
integral whose class depends on anything but the density topology, or on a Slater-geminal
integral whose class does not change at an equal-mass value of its exponent (guards
\texttt{tests/test\_paper59\_h2\_geminal\_master.py},
\texttt{tests/test\_paper63\_geminal\_shapes.py},
\texttt{tests/test\_paper63\_gy\_gaussian.py},
\texttt{tests/test\_paper63\_mg\_mechanism.py}).


With the dictionary of Sec.~\ref{sec:dictionary} the pattern can be stated
for every exponential-type orbital at once, and bounded.

\begin{conjecture}[the density-graph ladder]
\label{conj:ladder}
Let $\Gamma$ be a connected density graph over exponential-type orbitals with
Coulomb or Slater-geminal kernels, and let its fibre be the integral with
all centers coincident.  Then:\ (i) if $\Gamma$ is a tree of $n$ two-center
densities, the minimal annihilator of the fibre's one-scale family has, among
its composition factors, the operator of $F_n$ of
Sec.~\ref{sec:four_density}---the Hadamard product of the Legendre operator
with the equal-mass banana with $n-1$ propagators---and nothing above it;\
(ii) the rung names that operator, not a finite basis of periods:\ the fibre
is in general an iterated integral of its top piece against elementary
kernels, and no finite set of periods of the top piece and the rungs below
it is expected to contain the value;\ (iii) for a Coulomb ring, each loop
that integrates out exactly by uniqueness lowers the rung by one density, and
a geminal loop does not;\ (iv) the integrated value over separations and
Feynman parameters has class at most that of its fibre.
\end{conjecture}

\emph{What is proved, what is measured, what is open.}  Part (i) is proved
at the level of the holomorphic period for all $n$ (Eq.~\eqref{eq:tree} and
the Hadamard relation, Sec.~\ref{sec:four_density}), for every exponential
orbital by Sec.~\ref{sec:dictionary}, and for the real three-dimensional
integral at $n=3$ and $n=4$;\ for $n\ge5$ the three-dimensional reduction
is not shown.  Part (ii) is the lesson of the three-electron chain
(Sec.~\ref{sec:three_electron}):\ at equal scales the fibre collapses onto
K3 periods and $\pi^3$, but at four generic scale triples integer-relation
searches at 320 to 400 digits in five bases built from the K3 and elliptic
periods, their derivatives, products, dilogarithms, trilogarithms and Tate
numbers find nothing while every planted relation is recovered, and on both
one-parameter slices computed the fibre is provably a K3 function integrated
once against an elementary kernel.  That null is the conjecture's stated test
case:\ a relation found in a principled widening of the basis would show the
fibre sits in a finite period ring after all and part (ii) would be
withdrawn;\ a relation that needs a new period type would move the chain off
rung three.  The record is \texttt{docs/predictions/2026-10-19\_p3e2\_widened.md}.
Part (iii) rests on one ring, the triangle, in which the Coulomb loop drops
and the geminal loop does not (Sec.~\ref{sec:density_ring});\ the ring of
four was not computed, and the squashed ring's momentum dies as a power of
its aspect ratio.  Part (iv) holds in every instance:\ the two-center
exchange class builds the elliptic fibre and closes at genus zero, and the
integrated three-center value leaves the two-center class as a function of
the bond angle (Paper~59).  The conjecture fails on a Coulomb-only integral
whose class depends on anything but the density graph, on a tree whose
real three-dimensional reduction brings a piece above $F_n$, or on a
Slater-geminal integral whose class does not change at an equal-mass value
of its exponent.  The nearest theorem in shape remains
Brown's~\cite{brown2009periods}, for massless graphs;\ here the masses enter,
and a proof would have to run through the Symanzik polynomials of
$\Gamma$ with the Yukawa masses on its lines, which we have not attempted.

%% ========================================================================
\section{Placement in the literature}
\label{sec:placement}
%% ========================================================================

\textbf{[PRIOR ART]} The operators of the ladder are known.  The Hadamard
construction and the tables that place $F_4$ and $F_5$ at entries \#58 and
\#39 are Almkvist and Zudilin's~\cite{almkvist_zudilin2007,aesz2005};\ the
equal-mass banana operators are Vanhove's~\cite{vanhove2014} and their
holomorphic periods are Broadhurst's random-walk
moments~\cite{broadhurst2009walks};\ that an $l$-loop banana carries a
Calabi--Yau $(l-1)$-fold is B\"onisch, Fischbach, Klemm, Nega and
Safari's~\cite{bonisch2021};\ the two-scale and three-scale Bessel moments
are Bailey, Borwein, Broadhurst and Glasser's~\cite{bbbg2008}, and the
Kloosterman motives behind the Ising class are Fres\'an, Sabbah and
Yu's~\cite{fresansabbahyu2023}.  Integration by parts on the three-electron
tetrahedron is Puchalski, K\c{e}dziera and Pachucki's~\cite{pkp2009}, and
the kite's triangle threshold is Broadhurst's~\cite{broadhurst2022kites}.
What is new here is the route from molecular integrals to these objects, the
tree identity that makes the holomorphic period depend on $n$ alone, the
dictionary that makes the ladder independent of the orbitals, the readings
of the ring, the geminal and the Gaussian, and the prospective record.

\textbf{[PRIOR ART]} On the chemistry side the closest work is the Averys'
momentum-space Coulomb--Sturmian program, whose monograph~\cite{avery_avery_book2018}
and papers~\cite{avery2013,avery_avery2015,avery_avery_4center} treat
exactly these integrals under the same Fock projection.  Their route to a
multicenter repulsion integral is to expand each two-center density in a
one-center Coulomb--Sturmian basis of doubled exponent on its two centers,
solving a linear system whose matrix is the overlap of that basis;\ the
expansion is a series rather than an identity, its coefficient matrices are
not sparse, and at short internuclear separation the system is ill-conditioned
by the overcompleteness of the basis, which is handled by phasing out its
small eigenvalues~\cite{avery2013}.  In the language of this paper that route
reduces every integral to the genus-zero one-center class by an infinite
series, and the short-range difficulty it meets is where the series
confronts the genus jump of Eq.~\eqref{eq:curve}:\ an elliptic period is
not a finite combination of genus-zero terms.  Neither the Bessel-moment
reduction nor the ladder appears in that program;\ this is an absence from a
reading of the monograph's chapters on interelectron repulsion and of the
cited papers, not a proof of absence.  The overcompleteness itself is older
than either program~\cite{monkhorst_jeziorski1979}.

%% ========================================================================
\section{Scope, and what is open}
\label{sec:scope}
%% ========================================================================

Four caps bound every statement above.  First, the rungs are read on the
fibre and its holomorphic period;\ the integrated value may sit lower, and
in one class it does.  Second, the reduction of the real three-dimensional
integral to the Bessel-moment family is shown for the chain of three and the
star and path of four;\ for five or more densities the holomorphic-period
statement stands and the three-dimensional one is not shown.  Third, the
dictionary gives an upper bound on the rung of a general-orbital integral,
not its value;\ every computed case attains the bound, and none is a proof
that it must.  Fourth, among rings only the triangle was computed, and the
generic three-electron fibre's content beyond the K3 and elliptic periods is
evidence from integer-relation searches, not a theorem.

What is open is in two layers.  Below the conjecture:\ the three-dimensional
reduction for $n\ge5$;\ the ring of four;\ whether the rung of a
general-orbital integral is always attained;\ and the principled widening of
the P3e-2 basis.  Above it:\ a proof through the Symanzik polynomials that the
density graph with Yukawa masses has the motive the ladder measures.  The
finite closed form of the first rung's integrated value is Paper~59's open
problem and Paper~61's subject, not this paper's.

%% ========================================================================
\section{Conclusion}
%% ========================================================================

The class of a molecular integral over exponential-type orbitals is set by
its density graph.  Every such orbital is a derivative of the $1s$ Yukawa
function in its exponent and its center, so every integral is a derivative
of the $1s$ integral of its graph and cannot climb above it;\ the ladder is
orbital-independent.  On the graph, trees of $n$ linked two-center densities
share one holomorphic period, a Bessel moment whose one-scale operator is the
Legendre operator Hadamard the equal-mass banana with $n-1$ propagators:\
elliptic at two densities, a K3 at three, a Calabi--Yau threefold at four, and
nothing new at five.  The integrals of a molecular Hamiltonian link two
densities and stay elliptic at the fibre for every molecule;\ the class
climbs only in correlated integrals that link more.  A Coulomb ring of three
does not climb and a geminal ring does, read as uniqueness.  A Slater geminal
is an imaginary external momentum on this ladder and a Gaussian one leaves
it.  The operators are known;\ the route from chemistry to them, the tree
identity, the dictionary and the prospective record are what this paper adds,
and the conjecture that organizes them is stated with the case that tests
it:\ the rung names an operator, not a number, and the generic three-electron
fibre is where that must be decided.

%% ========================================================================
\begin{thebibliography}{99}

\bibitem{loutey_paper59}
J.~Loutey, ``The Three-Center Electron-Repulsion Integral Is an Elliptic
Bessel Moment:\ A Genus-One Radial Core at the Third Center in Momentum-Space
Slater Theory,'' GeoVac Paper~59 (2026).

\bibitem{loutey_paper57}
J.~Loutey, ``Forced and Free Content in the GeoVac Framework:\ An
Observations Paper,'' GeoVac Paper~57 (2026).

\bibitem{avery_avery_book2018}
J.~E.~Avery and J.~S.~Avery, \textit{Hyperspherical Harmonics and Their
Physical Applications} (World Scientific, Singapore, 2018).

\bibitem{monkhorst_jeziorski1979}
H.~J.~Monkhorst and B.~Jeziorski, ``No linear dependence or many-center
integral problems in molecular Sturmian calculations,'' \textit{J.\ Chem.\
Phys.}\ \textbf{71}, 5268--5269 (1979).


\bibitem{avery_avery_4center}
J.~E.~Avery and J.~S.~Avery, ``4-Center STO interelectron repulsion
integrals with Coulomb Sturmians,'' \textit{Adv.\ Quantum Chem.}\ \textbf{76},
133--146 (2018).

\bibitem{fromm_hill1987}
D.~M.~Fromm and R.~N.~Hill, ``Analytic evaluation of three-electron
integrals,'' \textit{Phys.\ Rev.\ A}\ \textbf{36}, 1013--1044 (1987).

\bibitem{remiddi1991}
E.~Remiddi, \textit{Phys.\ Rev.\ A}\ \textbf{44}, 5492 (1991).

\bibitem{chetyrkin_tkachov1981}
K.~G.~Chetyrkin and F.~V.~Tkachov, ``Integration by parts:\ the algorithm
to calculate $\beta$-functions in 4 loops,'' \textit{Nucl.\ Phys.\ B}
\textbf{192}, 159 (1981).

\bibitem{laporta2000}
S.~Laporta, ``High-precision calculation of multi-loop Feynman integrals by
difference equations,'' \textit{Int.\ J.\ Mod.\ Phys.\ A} \textbf{15}, 5087
(2000);\ arXiv:hep-ph/0102033.

\bibitem{chen2026fourier}
W.~Chen, L.~Silberman, U.~Segal and C.~Greif, ``Fourier transform-based
computation of two-electron four-center Coulomb repulsion integrals,''
arXiv:2610.12608 (2026).

\bibitem{broadhurst2022kites}
D.~Broadhurst, ``Multivariate elliptic kites and tetrahedral tadpoles,''
arXiv:2212.01962 (v2, 2024), to appear in the proceedings of LT-15 (Varna,
2023).

\bibitem{pkp2009}
M.~Puchalski, D.~K\c{e}dziera and K.~Pachucki, ``Ground state of Li and
Be$^+$ using explicitly correlated functions,'' \textit{Phys.\ Rev.\ A}
\textbf{80}, 032521 (2009);\ arXiv:0909.5542.

\bibitem{moriello2020}
F.~Moriello, ``Generalised power series expansions for the elliptic planar
families of Higgs + jet production at two loops,'' \textit{JHEP}
\textbf{01} (2020) 150;\ arXiv:1907.13234.

\bibitem{hidding2021}
M.~Hidding, ``DiffExp, a Mathematica package for computing Feynman integrals
in terms of one-dimensional series expansions,'' \textit{Comput.\ Phys.\
Commun.}\ \textbf{269}, 108125 (2021);\ arXiv:2006.05510.

\bibitem{liu_ma2023}
X.~Liu and Y.-Q.~Ma, ``AMFlow:\ a Mathematica package for Feynman integrals
computation via auxiliary mass flow,'' \textit{Comput.\ Phys.\ Commun.}\
\textbf{283}, 108565 (2023);\ arXiv:2201.11669.

\bibitem{hehre1969}
W.~J.~Hehre, R.~F.~Stewart and J.~A.~Pople, ``Self-consistent molecular-orbital methods.
I.\ Use of Gaussian expansions of Slater-type atomic orbitals,'' \textit{J.\ Chem.\ Phys.}\
\textbf{51}, 2657 (1969).

\bibitem{kutzelnigg1994}
W.~Kutzelnigg, ``Theory of the expansion of wave functions in a Gaussian basis,''
\textit{Int.\ J.\ Quantum Chem.}\ \textbf{51}, 447 (1994).

\bibitem{braess1995}
D.~Braess, ``Asymptotics for the approximation of wave functions by exponential sums,''
\textit{J.\ Approx.\ Theory} \textbf{83}, 93 (1995).

\bibitem{weniger2005}
E.~J.~Weniger, ``The spherical tensor gradient operator,'' arXiv:math-ph/0505018 (2005).

\bibitem{psi4_2020}
D.~G.~A.~Smith \textit{et al.}, ``PSI4 1.4:\ open-source software for
high-throughput quantum chemistry,'' \textit{J.\ Chem.\ Phys.}\
\textbf{152}, 184108 (2020);\ the fit is in
\texttt{psi4/src/psi4/libmints/integralparameters.cc}, class
\texttt{FittedSlaterCorrelationFactor}.

\bibitem{ppr2004}
K.~Pachucki, M.~Puchalski and E.~Remiddi, ``Recursion relations for
Hylleraas three-electron integral,'' \textit{Phys.\ Rev.\ A} \textbf{70},
032502 (2004);\ arXiv:physics/0405057.

\bibitem{bbbg2008}
D.~H.~Bailey, J.~M.~Borwein, D.~Broadhurst, and M.~L.~Glasser, ``Elliptic
integral evaluations of Bessel moments and applications,''
\textit{J.\ Phys.\ A}\ \textbf{41}, 205203 (2008); arXiv:0801.0891.

\bibitem{bailey1936}
W.~N.~Bailey, ``Some infinite integrals involving Bessel functions (II),''
J.\ London Math.\ Soc.\ \textbf{11}, 16--20 (1936).

\bibitem{brown2009periods}
F.~C.~S.~Brown, ``On the periods of some Feynman integrals,'' arXiv:0910.0114 (2009).

\bibitem{blochvanhove2015}
S.~Bloch and P.~Vanhove, ``The elliptic dilogarithm for the sunset graph,''
\textit{J.\ Number Theory}\ \textbf{148}, 328--364 (2015); arXiv:1309.5865.

\bibitem{fresansabbahyu2023}
J.~Fres\'an, C.~Sabbah, and J.-D.~Yu, ``Quadratic relations between Bessel
moments,'' \textit{Algebra Number Theory}\ \textbf{17}, 541--602 (2023);
arXiv:2006.02702.

\bibitem{avery2013}
J.~E.~Avery, ``Fast electron repulsion integrals for molecular Coulomb
Sturmians,'' \textit{Adv.\ Quantum Chem.}\ \textbf{67}, 129--151 (2013).

\bibitem{avery_avery2015}
J.~E.~Avery and J.~S.~Avery, ``Molecular integrals for exponential-type orbitals
using hyperspherical harmonics,'' \textit{Adv.\ Quantum Chem.}\ \textbf{70},
265--324 (2015).

\bibitem{aesz2005}
G.~Almkvist, C.~van Enckevort, D.~van Straten, and W.~Zudilin, ``Tables of
Calabi--Yau equations,'' arXiv:math/0507430 (2005).

\bibitem{cydb}
Calabi--Yau differential operator database, version 3.0,
\texttt{cydb.mathematik.uni-mainz.de}, entry 2.9 (AESZ \#58).

\bibitem{gegelia_vanstraten2024}
N.~Gegelia and D.~van Straten, ``Paramodular forms from Calabi--Yau
operators,'' arXiv:2408.10183 (2024), Tables~3 and~6.

\bibitem{matuzas2026}
J.~Matuzas, ``A Calabi--Yau period in a four-electron correlation integral:
modular physical fibre and a hypergeometric endpoint,'' Zenodo,
doi:10.5281/zenodo.23218203 (2026).

\bibitem{almkvist_zudilin2007}
G.~Almkvist and W.~Zudilin, ``Differential equations, mirror maps and zeta
values,'' in \textit{Mirror Symmetry V}, AMS/IP Studies in Advanced
Mathematics \textbf{38}, 481--515 (2007); arXiv:math/0402386.

\bibitem{vanhove2014}
P.~Vanhove, ``The physics and the mixed Hodge structure of Feynman
integrals,'' arXiv:1401.6438 (2014).

\bibitem{beukers_heckman1989}
F.~Beukers and G.~Heckman, ``Monodromy for the hypergeometric function
${}_nF_{n-1}$,'' \textit{Invent.\ Math.}\ \textbf{95}, 325--354 (1989).

\bibitem{lim_tu_yu2026}
Z.~L.~Lim, F.-T.~Tu and J.-D.~Yu, ``Galois representations and modularity
for twisted moments of Kloosterman sums,'' arXiv:2609.13576 (2026).

\bibitem{davydychev_delbourgo1998}
A.~I. Davydychev and R.~Delbourgo, ``A geometrical angle on Feynman
integrals,'' J.~Math.~Phys.~\textbf{39}, 4299 (1998), arXiv:hep-th/9709216.

\bibitem{nickel1978}
B.~G. Nickel, J.~Math.~Phys.~\textbf{19}, 542 (1978).

\bibitem{preti2018}
M.~Preti, ``STR:\ a Mathematica package for the method of uniqueness,''
arXiv:1811.04935 (2018).

\bibitem{broadhurst2009walks}
D.~Broadhurst, ``Bessel moments, random walks and Calabi--Yau equations,''
preprint (2009).

\bibitem{bourjaily2019}
J.~L.~Bourjaily, A.~J.~McLeod, M.~von Hippel, and M.~Wilhelm, ``Bounded
collection of Feynman integral Calabi--Yau geometries,'' \textit{Phys.\
Rev.\ Lett.}\ \textbf{122}, 031601 (2019); arXiv:1810.07689.

\bibitem{padhy2017}
B.~Padhy, ``Analytic evaluation of some three- and four-electron atomic
integrals involving $s$ STO's and exponential correlation with unlinked
$r_{ij}$'s,'' arXiv:1701.05701 (2017).

\bibitem{bonisch2021}
K.~B\"onisch, F.~Fischbach, A.~Klemm, C.~Nega, and R.~Safari, ``Analytic
structure of all loop banana integrals,'' \textit{J.\ High Energy Phys.}\
\textbf{05}, 066 (2021); arXiv:2008.10574.

\bibitem{loutey_paper18}
J.~Loutey, ``Spectral-Geometric Exchange Constants: Projecting Discrete
Spectra onto Continuous Manifolds,'' Paper~18 in this series.

\bibitem{loutey_paper34}
J.~Loutey, ``GeoVac as a Two-Layer Framework: Projection Taxonomy and the
Empirical Matches Catalog,'' Paper~34 in this series.

\bibitem{loutey_paper61}
J.~Loutey, ``The Bessel-Moment Period Algebra of the Three-Center Coulomb
Integral: Modular Structure and Cosmic-Galois Placement,'' Paper~61 in this
series.

\end{thebibliography}

\end{document}
